

%  For support, please attach all files needed for compiling as well as the log file, and specify your operating system, LaTeX version, and LaTeX editor.

%=================================================================
\documentclass[journal,article,submit,pdftex,moreauthors]{Definitions/mdpi} 
\usepackage[utf8]{inputenc}
\usepackage{amssymb}
\usepackage{algorithm}
\usepackage{algpseudocode}
\usepackage{threeparttable}
%\documentclass[preprints,article,submit,pdftex,moreauthors]{Definitions/mdpi} 
% For posting an early version of this manuscript as a preprint, you may use "preprints" as the journal. Changing "submit" to "accept" before posting will remove line numbers.

% Below journals will use APA reference format:
% admsci, aieduc, behavsci, businesses, econometrics, economies, education, ejihpe, famsci, games, humans, ijcs, ijfs, jintelligence, journalmedia, jrfm, jsam, languages, peacestud, psycholint, publications, tourismhosp, youth

% Below journals will use Chicago reference format:
% arts, genealogy, histories, humanities, laws, literature, religions, risks, socsci

%----------
% Class Options:
%----------
%-----
% journal
%-----
% Choose between the following MDPI journals:
% accountaudit, acoustics, actuators, addictions, adhesives, admsci, adolescents, aerobiology, aerospace, agriculture, agriengineering, agrochemicals, agronomy, ai, aichem, aieduc, aieng, aimater, aimed, aipa, air, aisens, algorithms, allergies, alloys, amh, analog, analytica, analytics, anatomia, anesthres, animals, antibiotics, antibodies, antioxidants, applbiosci, appliedchem, appliedmath, appliedphys, applmech, applmicrobiol, applnano, applsci, aquacj, architecture, arm, arthropoda, arts, asc, asi, astronautics, astronomy, atmosphere, atoms, audiolres, automation, axioms, bacteria, batteries, bdcc, behavsci, beverages, biochem, bioengineering, biologics, biology, biomass, biomechanics, biomed, biomedicines, biomedinformatics, biomimetics, biomolecules, biophysica, bioresourbioprod, biosensors, biosphere, biotech, birds, blockchains, bloods, blsf, brainsci, breath, buildings, businesses, cancers, carbon, cardio, cardiogenetics, cardiovascmed, catalysts, cells, ceramics, challenges, chemengineering, chemistry, chemosensors, chemproc, children, chips, cimb, civileng, cleantechnol, climate, clinbioenerg, clinpract, clockssleep, cmd, cmtr, coasts, coatings, colloids, colorants, commodities, complexities, complications, compounds, computation, computers, condensedmatter, conservation, constrmater, cosmetics, covid, crops, cryo, cryptography, crystals, csmf, ctn, culture, curroncol, cyber, dairy, data, ddc, dentistry, dermato, dermatopathology, designs, devices, dhi, diabetology, diagnostics, dietetics, digital, disabilities, diseases, diversity, dna, drones, dynamics, earth, ebj, ecm, ecologies, econometrics, economies, edm, education, eesp, ejihpe, electricity, electrochem, electronicmat, electronics, encyclopedia, endocrines, energies, eng, engproc, entomology, entropy, environments, environremediat, epidemiologia, epigenomes, esa, est, famsci, fermentation, fibers, fintech, fire, fishes, fluids, foods, forecasting, forensicsci, forests, fossstud, foundations, fractalfract, fuels, future, futureinternet, futurepharmacol, futurephys, futuretransp, galaxies, games, gases, gastroent, gastrointestdisord, gastronomy, gels, genealogy, genes, geographies, geohazards, geomatics, geometry, geosciences, geotechnics, geriatrics, germs, glacies, grasses, green, greenhealth, gucdd, hardware, hazardousmatters, healthcare, hearts, hemato, hematolrep, hep, heritage, higheredu, highthroughput, histories, horticulturae, hospitals, humanities, humans, hydrobiology, hydrogen, hydrology, hydropower, hygiene, idr, iic, ijcs, ijem, ijerph, ijfs, ijgi, ijmd, ijms, ijns, ijom, ijpb, ijt, ijtm, ijtpp, ime, immuno, informatics, information, infrastructures, inorganics, insects, instruments, inventions, iot, j, jaestheticmed, jal, jcdd, jcm, jcp, jcrm, jcs, jcto, jdad, jdb, jdream, jemr, jeta, jfb, jfmk, jgbg, jgg, jimaging, jintelligence, jlpea, jmahp, jmmp, jmms, jmp, jmse, jne, jnt, jof, joi, joitmc, joma, jop, joptm, jor, journalmedia, jox, jpbi, jphytomed, jpm, jrfm, jsam, jsan, jtaer, jvd, jzbg, kidneydial, kinasesphosphatases, knowledge, labmed, laboratories, lae, land, languages, laws, life, lights, limnolrev, lipidology, liquids, literature, livers, logics, logistics, lubricants, lymphatics, machines, macromol, magnetism, magnetochemistry, make, marinedrugs, materials, materproc, mathematics, mca, measurements, medicina, medicines, medsci, membranes, merits, metabolites, metals, meteorology, methane, metrics, metrology, micro, microarrays, microbiolres, microelectronics, micromachines, microorganisms, microplastics, microwave, minerals, mining, mmphys, modelling, molbank, molecules, mps, msf, mti, multimedia, muscles, nanoenergyadv, nanomanufacturing, nanomaterials, ncrna, ndt, network, neuroglia, neuroimaging, neurolint, neurosci, nitrogen, notspecified, nursrep, nutraceuticals, nutrients, obesities, occuphealth, oceans, ohbm, onco, optics, oral, organics, organoids, osteology, oxygen, pandemics, parasites, parasitologia, particles, pathogens, pathophysiology, peacestud, pediatrrep, pets, pharmaceuticals, pharmaceutics, pharmacoepidemiology, pharmacy, philosophies, photochem, photonics, phycology, physchem, physics, physiologia, plants, plasma, platforms, pollutants, polymers, polysaccharides, populations, poultry, powders, precipitation, precisoncol, preprints, proceedings, processes, prosthesis, proteomes, psf, psychiatryint, psychoactives, psycholint, publications, purification, quantumrep, quaternary, qubs, radiation, rdt, reactions, realestate, receptors, recycling, regeneration, religions, remotesensing, reports, reprodmed, resources, rheumato, risks, rjpm, robotics, rsee, ruminants, safety, sci, scipharm, sclerosis, seeds, sensors, separations, sexes, shi, signals, sinusitis, siuj, skins, smartcities, sna, societies, socsci, software, soilsystems, solar, solids, spectroscj, sports, standards, stats, std, stratsediment, stresses, surfaces, surgeries, suschem, sustainability, symmetry, synbio, systems, tae, targets, taxonomy, technologies, telecom, test, textiles, thalassrep, therapeutics, thermo, timespace, tomography, tourismhosp, toxics, toxins, tph, transplantology, transportation, traumacare, traumas, tri, tropicalmed, universe, urbansci, uro, vaccines, vehicles, venereology, vetsci, vibration, virtualworlds, viruses, vision, waste, water, welding, wem, wevj, wild, wind, women, world, youth, zoonoticdis

%----
% article
%----
% The default type of manuscript is "article", but can be replaced by: 
% abstract, addendum, article, benchmark, book, bookreview, briefcommunication, briefreport, casereport, changes, clinicopathologicalchallenge, comment, commentary, communication, conceptpaper, conferenceproceedings, correction, conferencereport, creative, datadescriptor, discussion, entry, expressionofconcern, extendedabstract, editorial, essay, erratum, fieldguide, hypothesis, interestingimages, letter, meetingreport, monograph, newbookreceived, obituary, opinion, proceedingpaper, projectreport, reply, retraction, review, perspective, protocol, shortnote, studyprotocol, supfile, systematicreview, technicalnote, viewpoint, guidelines, registeredreport, tutorial,  giantsinurology, urologyaroundtheworld
% supfile = supplementary materials

%-----
% submit
%-----
% The class option "submit" will be changed to "accept" by the Editorial Office when the paper is accepted. This will only make changes to the frontpage (e.g., the logo of the journal will get visible), the headings, and the copyright information. Also, line numbering will be removed. Journal info and pagination for accepted papers will also be assigned by the Editorial Office.

%--------
% moreauthors
%--------
% If there is only one author the class option oneauthor should be used. Otherwise use the class option moreauthors.

%----
% pdftex
%----
% The option pdftex is for use with pdfLaTeX. Remove "pdftex" for (1) compiling with LaTeX & dvi2pdf (if eps figures are used) or for (2) compiling with XeLaTeX.

%=================================================================
% MDPI internal commands - do not modify
\firstpage{1} 
\makeatletter 
\setcounter{page}{\@firstpage} 
\makeatother
\pubvolume{1}
\issuenum{1}
\articlenumber{0}
\pubyear{2026}
\copyrightyear{2026}
%\externaleditor{Firstname Lastname} % More than 1 editor, please add `` and '' before the last editor name
\datereceived{ } 
\daterevised{ } % Comment out if no revised date
\dateaccepted{ } 
\datepublished{ } 
%\datecorrected{} % For corrected papers: "Corrected: XXX" date in the original paper.
%\dateretracted{} % For retracted papers: "Retracted: XXX" date in the original paper.
%\doinum{} % Used for some special journals, like molbank
%\pdfoutput=1 % Uncommented for upload to arXiv.org
%\CorrStatement{yes}  % For updates
%\longauthorlist{yes} % For many authors that exceed the left citation part
%\IsAssociation{yes} % For association journals

%=================================================================
% Add packages and commands here. The following packages are loaded in our class file: fontenc, inputenc, calc, indentfirst, fancyhdr, graphicx, epstopdf, lastpage, ifthen, float, amsmath, amssymb, lineno, setspace, enumitem, mathpazo, booktabs, titlesec, etoolbox, tabto, xcolor, colortbl, soul, multirow, microtype, tikz, totcount, changepage, attrib, upgreek, array, tabularx, pbox, ragged2e, tocloft, marginnote, marginfix, enotez, amsthm, natbib, hyperref, cleveref, scrextend, url, geometry, newfloat, caption, draftwatermark, seqsplit
% cleveref: load \crefname definitions after \begin{document}

\providecommand{\cmark}{$\bullet$}
\providecommand{\omark}{$\circ$}
\providecommand{\xmark}{$\times$}

%=================================================================
% Please use the following mathematics environments: Theorem, Lemma, Corollary, Proposition, Characterization, Property, Problem, Example, ExamplesandDefinitions, Hypothesis, Remark, Definition, Notation, Assumption
%% For proofs, please use the proof environment (the amsthm package is loaded by the MDPI class).

%=================================================================
% Full title of the paper (Capitalized)
\Title{LACE: Latent Commonality Expectation-Maximization for Box-supervised Tree Crown Instance Segmentation}

% Author Orchid ID: enter ID or remove command
\newcommand{\orcidauthorA}{0000-0000-0000-000X} % Add \orcidA{} behind the author's name
%\newcommand{\orcidauthorB}{0000-0000-0000-000X} % Add \orcidB{} behind the author's name

% Authors, for the paper (add full first names)
\Author{Thomas Pitts $^{1}$\orcidA{}, Kunqi Li $^{2}$ and Bin Liang $^{2,}$*}

%\longauthorlist{yes}

% MDPI internal command: Authors, for metadata in PDF
\AuthorNames{Thomas Pitts, Kunqi Li and Bin Liang}

%\longauthorlist{yes}

% Affiliations / Addresses (Add [1] after \address if there is only one affiliation.)
\address[1]{%
$^{1}$ \quad Department of Data Science, University of Technology Sydney; thomas.pitts@uts.edu.au\\
$^{2}$ \quad Affiliation 2; e-mail@e-mail.com}

% Contact information of the corresponding author
\corres{Correspondence: e-mail@e-mail.com; Tel.: (optional; include country code; if there are multiple corresponding authors, add author initials) +xx-xxxx-xxx-xxxx (F.L.)}

% Current address and/or shared authorship
%\firstnote{Current address: Affiliation.}  
% Current address should not be the same as any items in the Affiliation section.

%\secondnote{These authors contributed equally to this work.}
% The commands \thirdnote{} till \eighthnote{} are available for further notes.

%\simplesumm{} % Simple summary

%\conference{} % An extended version of a conference paper

% Abstract (Do not insert blank lines, i.e. \\) 
\abstract{Individual tree crown segmentation underpins key ecological applications globally, including forest management and carbon accounting. However, existing models are predominantly trained on dense canopy forest imagery and degrade in savanna and drylands, where tree crowns are sparse, of variable appearance, and underrepresented in annotated benchmarks. These models also depend on costly polygon crown annotations. We introduce LACE (LAtent Commonality Expectation-maximization), a box-supervised instance segmentation model for 0.1 m/px aerial RGB imagery. LACE uses a frozen DINOv3-web ViT-L/16 encoder, applied at four spatial offsets and interlaced into a denser feature grid, with a lightweight CenterNet-style detection head trained only on bounding boxes. Masks are recovered by a latent commonality box-to-mask module: foreground and background mixtures are estimated in whitened feature space using expectation-maximization, separating recurring appearance within bounding boxes from its surroundings, without tree- or domain-specific priors. On Restor Foundation's OAM-TCD holdout, LACE reaches a mask mAP50 of $0.630 \pm 0.005$ (3 seeds) trained on 900 images and no mask annotations, above the 0.586 scored by Restor's released mask-supervised Mask R-CNN, which was trained on the full $\sim$4.2k image set (4169 images). On NEONTreeEvaluation, using the official evaluation code, LACE reaches $0.728 \pm 0.003$ F1@0.4 (5 seeds) from RGB boxes alone, exceeding the authors' DeepForest model's published 0.719 while using $\sim$10\% (TBC) of the annotations and none of the LiDAR or hyperspectral channels. By matching or surpassing fully-supervised specialist baselines from boxes alone, LACE offers annotation-efficient tree crown instance segmentation in dryland biomes where labeled data is scarce.}

% Keywords
\keyword{keyword 1; keyword 2; keyword 3 (List three to ten pertinent keywords specific to the article; yet reasonably common within the subject discipline.)} 

% The fields PACS, MSC, and JEL may be left empty or commented out if not applicable
%\PACS{J0101}
%\MSC{}
%\JEL{}

%%%%%%%%%%%%%%%%%%%%%%%%%%%%%%%%%%%%%%%%%%
% Only for the journal Diversity
%\LSID{\url{http://}}

%%%%%%%%%%%%%%%%%%%%%%%%%%%%%%%%%%%%%%%%%%
% Only for the journal Applied Sciences
%\featuredapplication{Authors are encouraged to provide a concise description of the specific application or a potential application of the work. This section is not mandatory.}
%%%%%%%%%%%%%%%%%%%%%%%%%%%%%%%%%%%%%%%%%%

%%%%%%%%%%%%%%%%%%%%%%%%%%%%%%%%%%%%%%%%%%
% Only for the journal Data
%\dataset{DOI number or link to the deposited data set if the data set is published separately. If the data set shall be published as a supplement to this paper, this field will be filled by the journal editors. In this case, please submit the data set as a supplement.}
%\datasetlicense{License under which the data set is made available (CC0, CC-BY, CC-BY-SA, CC-BY-NC, etc.)}

%%%%%%%%%%%%%%%%%%%%%%%%%%%%%%%%%%%%%%%%%%
% Only for the journal BioTech, Fishes, Neuroimaging and Toxins
%\keycontribution{The breakthroughs or highlights of the manuscript. Authors can write one or two sentences to describe the most important part of the paper.}

%%%%%%%%%%%%%%%%%%%%%%%%%%%%%%%%%%%%%%%%%%
% Only for the journal Encyclopedia
%\encyclopediadef{For entry manuscripts only: please provide a brief overview of the entry title instead of an abstract.}

%%%%%%%%%%%%%%%%%%%%%%%%%%%%%%%%%%%%%%%%%%
% Different journals have different requirements. Please check the specific journal guidelines in the "Instructions for Authors" on the journal's official website.
%\addhighlights{yes}
%\renewcommand{\addhighlights}{%
%
%\noindent The goal is to increase the discoverability and readability of the article via search engines and other scholars. Highlights should not be a copy of the abstract, but a simple text allowing the reader to quickly and simplified find out what the article is about and what can be cited from it. Each of these parts should be devoted up to 2~bullet points.\vspace{3pt}\\
%\textbf{What are the main findings?}
% \begin{itemize}[labelsep=2.5mm,topsep=-3pt]
% \item First bullet.
% \item Second bullet.
% \end{itemize}\vspace{3pt}
%\textbf{What are the implications of the main findings?}
% \begin{itemize}[labelsep=2.5mm,topsep=-3pt]
% \item First bullet.
% \item Second bullet.
% \end{itemize}
%}

%%%%%%%%%%%%%%%%%%%%%%%%%%%%%%%%%%%%%%%%%%
\begin{document}

\section{Introduction}

%The introduction should briefly place the study in a broad context and highlight why it is important. It should define the purpose of the work and its significance. The current state of the research field should be reviewed carefully and key publications cited. Please highlight controversial and diverging hypotheses when necessary. Finally, briefly mention the main aim of the work and highlight the principal conclusions. As far as possible, please keep the introduction comprehensible to scientists outside your particular field of research. Citing a journal paper \citep{ref-journal}.  Now citing a book reference \citep{ref-book1,ref-book2} or other reference types \citep{ref-unpublish,ref-url}. Please use the command \citep{ref-proceeding,ref-thesis} for the following MDPI journals, which use author--date citation: Administrative Sciences, Arts, Behavioral Sciences, Businesses, Econometrics, Economies, Education Sciences, European Journal of Investigation in Health, Psychology and Education, Games, Genealogy, Histories, Humanities, Humans, IJFS, Journal of Intelligence, Journalism and Media, JRFM, Languages, Laws, Literature, Psychology International, Publications, Religions, Risks, Social Sciences, Tourism and Hospitality, Youth. 

%%%%%%%%%%%%%%%%%%%%%%%%%%%%%%%%%%%%%%%%%%

Accurate measurement of tree crown cover can serve as a proxy for woody biomass estimation via allometries (Jucker et al., 2022) and thereby enable monitoring of carbon sequestration within areas containing trees since cabon content can be derived from woody biomass. Measurement techniques in this domain have traditionally focussed on dense canopy areas and forests, but there is an increasing body of literature that highlights the limitations of this approach as distribution, density, cover, and carbon content of trees in sparse-canopy biomes such as savanna and dryland are not well understood at sub-continental to continental scales (Tucker 2023 1-14), and in the Global South as a whole (Baudchon SELVABox). Recent studies highlight that the global contribution of sparse-canopy as a carbon sink are underestimated due to the underestimation of tree mass in sparse-canopy areas (Tucker 2023, Brandt 2016 + 2020, Crowther 2015, Bastin 2017). Moreover, approximately 1/3 of all trees globally are outside forests i.e. individual \& sparsely distributed at <10\% tree canopy cover (cite). 

 Several conventional methodologies have been applied to the measurement of tree cover in arid and semi-arid dryland areas, including local field surveys (Brandt 2020), Area sampling using (Reddy 2011, Raty 2020), and hybrid approaches combining field plots and remote sensing metrics (Mayamanikandan when?). Machine learning and deep learning applied to satellite and UAV data have since shifted the field from aggregate cover fractions toward mapping every individual tree: Brandt et al. (2020) mapped over 1.8 billion crowns across 1.3 million km² of the Sahara and Sahel from 0.5 m imagery; Tucker et al. (2023) extended this to carbon over 9.9 billion trees. Reiner et al. (2023) found that 29\% of Africa's tree cover lies outside previously classified forest. Mugabowindekwe et al. (2023) produced nation-wide tree-level carbon estimates for Rwanda. These studies establish that individual crown detection in open, scattered-tree drylands is tractable, but depend on proprietary sub-metre imagery and bespoke labels, motivating more transferable and label-efficient methods.

Modern deep learning approaches focus on instance segmentation on RGB imagery: DeepForest (Weinstein et al., 2019, 2020) established a RetinaNet-based baseline evaluated on the NeonTreeEvaluation benchmark (Weinstein et al., 2021), while Detectree2 (Ball et al., 2023) applied Mask R-CNN (He et al., 2017) to tropical crown polygons. Tong \& Zhang (2025) applied StarDist (cite original authors), a U-Net (cite) backbone model applying star-convex polygons as proposals during the detection stage rather than the more commonly used axis-aligned bounding boxes, allowing for improved Non-Maximum Suppression (NMS) performance and more accurate detections. This approach was designed as a solution to the challenging problem of individual tree crown instance segementation in dense forest canopy, where crowns are often touching and overlapping, which can contribute to many valid detections being needlessly dropped due to the coarse tessellation of bounding box rectangles. Similarly, the authors of the tropical rainforest SelvaMask dataset (Duguay et al., 2026) utilized a DINO (DIstillation with NO labels - Caron et al., 2021) based tree crown detection model coupled with a Segment Anything Model (SAM) (Meta cite) for prediction mask generation, applied to instance segmentation in dense canopy rainforests. This approach using SAM has also been explored in a number of works (RSPrompter, Chen et al., 2025, and Teng et al. (cite x2 workshop ICLR paper and BalSAM) both with and without the integration of a Digital Surface Model (DSM) depth map, who reported suboptimal results for out-of-the-box SAM but found further finetuning promising.

We aim to address the tree crown instance segmentation space from two angles underexplored in the literature, both largely stemming from the lack of availability of data. Firstly, existing approaches focus on dense-canopy forests, and overwhelmingly exclude savanna and dryland biomes. For example, the authors of Detectree2 deliberately exclude images with <= 40\% tree canopy cover from their datasets (check exact quote and cite). In the DeepForest paper, Weinstein et al. highlight Onaqui, Utah (ONAQ) as the model's worst performing site (Weinstein et al., 2020) since it is comprised of "desert scrub site with a different vegetation structure from any of the training data". Secondly, in order to produce output individual tree crown instance polygon masks, existing models typically rely on mask annotations for training (cite 5 models here and/or literature reviews), which need to be produced manually by human annotators.


\subsection{Lack of sparse-cover environments}

We recognise a number of contributing factors to the exclusion of sparse-cover biomes:
\begin{enumerate}
\item	the limitations of spatial resolution, since an accurate model of ITC cover is hard to create with most publicly available satellite RGB datasets at 5-10 m/px. High (<1m/px) or ultra high (<= 0.1m/px) resolution is required (Bastin 2017 (check where)). For reference, in the OAM-TCD holdout set used in this work (439 images, 25,705 individual tree crowns), downsampling from the images' native 0.1m/px to a Ground Sample Distance (GSD) of 1m (i.e. the real-world area corresponding to one pixel) results in 24.2\% of tree crowns having zero pixels enclosed inside the crown boundary. At 2m GSD, 53.4\% of tree crowns enclose zero pixels, and 22.6\% of crowns are too small to produce even a single pixel mask;
\item	underrepresentation in high quality annotated datasets - both due to a) lack of focus in the remote sensing community generally, and furthermore b) in the datasets that do exist, fewer samples by definition relative to open and closed forest canopy (UN FAO define 10-39\% tree canopy cover as 'open forest' and >=40\% as 'closed forest'). On a study of 9.9 billion individual tree crowns in the Sahara/Sahel/Sudan region, Tucker et al (2023) found that "areas with scattered trees are often represented by zero values" (cite Nature 615 p.85), attributing this to "the fact that previous models are rarely developed, trained and validated with plots of very sparse tree cover". The work of Tucker et al. (other papers to cite?) has been crucial in highlighting this vicious cycle: existing models trained on dense canopy underrepresent the amount of tree cover in sparse environments (Avitabile, Baccini, Bouvet (Tucker ref 6-8)), which in turn stymies further exploration of these biomes for tree crown detection and instance segmentation.
\end{enumerate}
%insert table showing datasets and biomes, as well as bounding boxes or masks - VERY SIMILAR TO SELVAMASK table

In short, the development of Remote Sensing Computer Vision techniques in dryland areas for the purposes of Tree Canopy monitoring has been underexplored, and yet represents a significant biome in measurement of global carbon sequestration. Recent annotated datasets such as SelvaBox (Baudchon et al., 2026) \& SelvaMask (Duguay et al., 2026) have broadened geographic coverage, and Restor Foundation's OAM-TCD dataset (Veitch-Michaelis et al., 2024) in particular incorporates a variety of biomes including urban, tundra, and savannah.

\subsection{Reliance on manual mask annotations}

% 1) Existing models focus on forests, neglect savanna & dryland -> 2) all rely on mask annotations which are expensive and slow. We present a model that only needs box annotations to train and uses the embedding space signature of the box annotation contents to automatically delineate individual tree crowns without explicit demonstration, and test this over global datasets including forest, savanna/dryland, and urban biomes. Leveraging embedding space vectors allows model to perform well in noisy and unseen contexts.


Tree crown instance segmentation also remains understudied due to the lack of individual tree crown mask annotations in datasets (Teng et al. 2025). Since crown polygons are slower and typically more expensive to annotate than boxes, box-supervised instance segmentation (Tian et al., 2021; Li et al., 2024; Cheng et al., 2023) offers a practical route to instance-level delineation without polygon labels. The recent rise of self-supervised backbones such as DINOv2 (Oquab et al., 2023) and DINOv3 (Siméoni et al., 2025) have demonstrated that models pretrained on supersized sets of diverse web imagery (1.7 billion images in the case of DINOv3-web) match or exceed Earth-observation-specific foundation models such as DINOv3-sat on high-resolution RGB tasks, with the latter retaining an advantage chiefly where multispectral or multitemporal signal is essential (Lacoste et al., 2023). Apples-to-apples comparison across this literature remains difficult however, as papers variously report mAP50, mAP50–95, instance- or semantic-F1 at differing IoU thresholds against boxes and/or polygons.


In this work we propose a 3-stage model LACE built on a frozen DINOv3-web (ViT-L/16px) (cite) encoder, with a 4m parameter CenterNet-style (cite) detector head, and a novel lightweight Expectation-maximization module that leverages latent commonality of the ground truth bounding box annotations. In other words, we ask the question "can we leverage the vector space signature of the tree-ness within DINO latent feature vectors, based only on the commonality between rectangular bounding box examples, to isolate the trees/foreground from the background?" Our experiments across a range of biomes including temperate, tropical, xeric, and urban scenes, demonstrate that even with diverse tree species and background context, a latent commonality module is sufficient to infer and delineate objects of interest from rough rectangular annotations alone. 

LACE achieves a mask mAP50 of $0.630 \pm 0.005$ on the OAM-TCD holdout benchmark and a box F1 of $0.728 \pm 0.003$ on NEONTreeEvaluation (evaluated at IOU=0.4) with our LACE model. The creators of NEONTreeEvaluation propose 0.4 IOU as a suitable threshold for tree crown detection tasks based on comparative analysis of agreement and overlap between different human annotations of the same datasets, and we follow their convention on this dataset for direct comparison with their results. 

By using only the hand-annotated bounding box annotations (which are sampled from 13 of the total 22 sites present in the test dataset), our model was able to match and exceed DeepForest's published results of 0.719 F1, evaluated at IOU>=0.4. We also did not make use of the a ‘canopy height model’ (CHM) (a LiDAR derived height raster at 1m spatial resolution), or any of the NEONTreeEvaluation dataset's hyperspectral imagery, since our focus is the underexplored and data-scarce sparse-canopy context. By demonstrating the efficacy of LACE on RGB data alone, we hope our contribution to the remote sensing community will have a broader application to data-scarce contexts. 

In summary, the lack of publicly available tree crown-annotated datasets generates a need for weakly supervised/low-annotation, performant TCIS models suitable for diverse open/sparse canopy contexts. Our contribution is deliberately orthogonal to dense-canopy tropical forest models, which we use as baselines for comparison in the Results section. In broad terms, most existing state of the art models focus on improving the precision of delineation in dense environments - by contrast, our focus is on identifying tree crowns in open landscapes. Nevertheless, LACE is still able to demonstrate SOTA performance in apples-to-apples comparisons, even with less training data and box supervision. By leveraging the rich DINOv3 features which enable rich cross-biome transferability, LACE is, to the best of our knowledge, the strongest performing biome-agnostic model in the TCIS space as evidenced by results on the OAM-TCD benchmark.


In this work, we present:

\begin{itemize}
    \item a) our model LACE, built on a DINOv3-web (ViT-L/16) backbone (cite), with a trained CenterNet-style detector head (Zhou et al., 2019) that identifies crown centers for detection based on the DINO output embeddings. We assess the performance of this approach compared to existing on two publicly available datasets Restor's OAM-TCD and DeepForest's NEON annotation dataset. We also share results of fine-tuning a DINOv3-web encoder on unseen unlabelled data, and find no improvement (nor regression) in doing so, adding a confirmatory result to DINO's authors' suggestion that 1) the model is essentially saturated at cross-domain tasks and 2) that the Remote Sensing-focussed DINOv3-sat actually underperforms on downstream tasks of TCD and TCIS in high resolution aerial RGB data.

    \item b) a novel box-to-mask module that is designed to extract the commonality in latent space using an Expectation-Maximization (EM) algorithm (cite) to extract 'treeness' within the bounding box. Unlike EM-Adapt (Papandreou et al., 2015), which amortizes the E-step into a gradient-trained DCNN with implicit cross-box sharing, we fit an explicit closed-form commonality mixture over frozen features. We exploit the similarities between signatures of composite embeddings within ground truth bounding boxes, following previous works such as DiscoBox (cite), DDT (cite), and STEGO (cite). In our model however, we utilise a World Model/JEPA (cite) style approach by eschewing pixel-space calculations in favour of strictly latent space analysis. This approach is made possible by the capabilities of Meta's DINOv3 encoder foundation model (cite), a Vision-transformer-based model trained on 1.7 Billion images using Self-supervised Learning.
\end{itemize}

\section{Related Work}

The use of computer vision models to process aerial and satellite RGB imagery for the purposes of tree crown measurements has seen a number of advancements in recent years (Zhao 2023). Two important benefits of RGB methods above other data modalities such as LIDAR or hyperspectral data-based methods are firstly the widespread availability of aerial RGB imagery, and secondly the high spatial resolution of that data up to and including submeter per pixel precision. Moreover, aerial RGB data can also be collected using UAV such as drones with relative ease. Consequently, in recent years there has been a focus in the RS community on RGB datasets which can directly leverage Computer Vision deep learning models (Zhu 2017, Zhao 2023) (SelvaBox refs: Weinstein 2021 NeonTreeEval) especially in the use of Convolutional Neural Networks (CNNs) and Vision Transformers (ViT, cite). 

\paragraph{Note on terminology}

We follow Ball et al. 2023 in defining Instance Segmentation as the precise delineation of individual tree crowns. Since individual trees are the fundamental unit that aggregate to the higher-scale biomass (Yang et al. 2022b, Fu et al. 2024), regardless of biome, we focus on Tree Crown Instance Segmentation (TCIS). Other related categories include Tree Crown Detection (TCD) models, which aim to detect instances of a given object and delineate them with a rectangular bounding box e.g. YOLO family (cite), and Semantic Segmentation (TCS) models which aim to categorise each pixel into a class. 

\subsection{Tree Crown Models}

\paragraph{Convolutional Neural Networks}

Modern CNN-based object detection models are comprised of three parts: 
\begin{enumerate}
\item a feature extractor network ('backbone') which is a neural network that scans with convolutions to pick up on features in the image, with simpler features typically extracted earlier in the network's layers and more complex composite features requiring greater depth (Goldblum 2023);
\item a 'neck', which can merge the detections across different layers of the backbone and provide access to the whole gamut of feature scale and complexity;
\item the Detector 'head(s)', which take the encoded features as inputs and produce the entire CNN's output, which can be, for example, bounding boxes around detected objects, polygons for a precise segmentation mask, or classification categories. Detector heads are generally either one-stage or two-stage detectors. Two-stage detectors first utilise a module to create many possible bounding boxes which are region proposals, before the second stage module extracts features from the proposals (Soviany, 2018). In this way, the model can accurately detect within a large image by focussing on high probability smaller regions. 
\end{enumerate}


Within TCD, widely used two-stage detector CNNs include the Region-based CNN (R-CNN) family of models (Girshick et al., 2014) - current best in class models include Fast R-CNN, Faster R-CNN and Mask R-CNN. One-stage detectors include YOLO ('You Only Look Once') series of models (Redmon; Bochkovskiy),  RetinaNet (Lin 2017) and EfficientDet (Tan), DetectNet (cite). 

\subsubsection{Star-convex}

Other significant contributions in the field include (Tong \& Zhang, 2025), based on StarDist (Schmidt et al., 2018), which demonstrated the TCIS improvements provided by useage of star-convex polygons during Non-Maximum Suppression (NMS) stage whereby initial predictions with significant overlap are pruned, eliminating lower confidence predictions through a variety of methods and strategies. Tong \& Zhang were able to show benefits in dense canopy TCIS on RGB images relative to Mask R-CNN and other SOTA models, specifically for overlapping dense-canopy tree crowns. MP-PolarMask (2024 cite) extends this idea to use multiple crown centers which better facilitates tree crowns that are not roughly star-convex shaped e.g. palm trees.  ZS-TreeSeg/FG-TreeSeg (Chen et al., 2026) also models tree crowns as star-convex objects, and introduces a vector field flow-based crown separation mechanism to improve instance segmentation of touching tree crowns and dense canopy.

\paragraph{DeepForest}

DeepForest (Weinstein et al., cite) is a significant baseline model in Tree Crown Detection (TCD). DeepForest is a RetinaNet on a ResNet50 backbone and Feature Pyramid Network (FPN) trained on over 30m algorithmically generated crowns from 22 National Ecological Observatory Network (NEON) sites and further fine-tuned using 10,000 hand-annotated crowns. % more detail on DF including Hyperspectral, LIDAR training data

\subsubsection{Pretrained Foundation Models}

Detectron2 is the %TBC


Meta's Segment Anything Model (SAM, cite), now in its third iteration (cite) %TBC


Meta's latest version of the Self-Distillation with No Labels (DINO) family is DINOv3 (Siméoni et al., 2025): a SOTA encoder trained with fully Self-Supervised Learning on 1.7 billion images pulled from various online sources. Utilizing SSL faciliotates training regimes on orders of magnitude more image data since weakly and fully supervized methods (Radford et al., 2021; Dehghani et al., 2023; Bolya et al., 2025) require high quality labels and/or metadata for the training images. Highly context-agnostic, DINOv3-web is an example of a single frozen SSL backbone that can serve as a "universal visual encoder" (Siméoni et al., 2025) capable of generating rich embedding vectors that capture the semantic information from the raw pixel patches in context. By incorporating both a global image-level objective and a local iBOT-style (Zhou et al., 2021) patch-level latent reconstruction objective in the loss function, results in a model that excels at encoding global and local features. In this work, we use the DINOv3-web ViT-L condensed model, with native 16px patch. 


% run DeepForest on TCD439? Compare box AP40? Also SelvaBox? Best SelvaBox BOX detection map50-95 is 44.29(±0.33) but they remove 0-annotation tiles, and mask canopy entirely. 

- SELVABOX (Baudchon 2026) is a UAV-captured RGB dataset of tropical dense-canopy forest in Panama, Brazil, and Ecuador, composed of 83,000 ITC manual bounding box annotations. Its authors publish a DINO-Swin-L based model trained onNeonTreeEvaluation (see Datasets), OAM-TCD, and another dataset called QuebecTrees, which achieves a best result of box mAP50-95 of $44.29 \pm 0.33$ on OAM-TCD holdout set, improving upon DeepForest ($39.00 \pm 0.21$) and a Faster R-CNN ResNet50 ($38.34 \pm 0.26$) finetuned on the same data. A second related dataset, SELVAMASK (Duguay, Baudchon et al., 2026) provides ITC instance segmentation masks and assesses the SelvaBox detection model (also referred to by the authors as 'SelvaBox') in combination with both frozen \& fine-tuned SAM3 decoder module on the OAM-TCD dataset for TCIS task. The authors compare results on a number of benchmark datasets including OAM-TCD alongside Detectree2 and DeepForest with a similar frozen SAM3 mask prediction module. SELVAMASK's focus is increased precision of tree crown boundaries in dense canopy, whereas our focus is on open canopy detection accuracy. They note that their model is specialized "to dense, interlocking canopies, slightly reducing its transferability to temperate forests". (Check the GSD transferability to OAM-TCD? SelvaMask dataset is 1.3-3.5 cm/px whereas OAM-TCD is 10cm/px). 


%Other methods enhance SAM backbones via learned prompts with RSPrompter (Chen et al., 2024), leveraging digital surface models with BalSAM (Teng et al., 2025) or use LiDAR point clouds as point-prompts (Zhu et al., 2025). - SelvaMask p.3

Leveraging SAM as an out-of-the-box instance segmenter for remote sensing has been explored in RSPrompter (Chen et al., 2023), who overcome SAM's reliance on prompts such as points/boxes/masks by training a prompter module that learns the intermediate layer features of the SAM encoder on labelled training images, and then using those embeddings as input to SAM's mask decoder in order to produce category-labelled instance masks. 

- 
- BalSAM (Teng et al (Bringing SAM to new heights)) Teng also finds poor zero-shot SAM performance
- FM-SAM (Que 2026) 
- Tree-SAM (Huang 2026) (urban and mixed-biome!)
- Zeroshot (Chen 2025)
- TreePseCo (Vaschetti 2025)
- Ginkgo (Mengyuan 2025)



\begin{table}[t]
\centering\footnotesize
\setlength{\tabcolsep}{3.5pt}
\renewcommand{\arraystretch}{1.15}
\caption{Supervision, inference cost, and evaluated biome coverage of
representative individual tree crown delineation methods.}
\label{tab:gap}
\begin{tabular}{@{}llllcccc@{}}
\toprule
 & & & 2nd & GSD & \multicolumn{3}{c}{Biome} \\
\cmidrule(l){6-8}
Method & Sup. & Input & net & (cm) & Tr & Te & Sv \\
\midrule
\multicolumn{8}{l}{\textit{Polygon-supervised, end-to-end}} \\
Mask R-CNN \cite{oamtcd} & Polygon & RGB & -- & 10 & \omark & \omark & \omark \\
detectree2 \cite{ball2023} & Polygon & RGB & -- & 10 & \cmark & \xmark & \xmark \\
Mask2Former \cite{dersch2023} & Polygon & MS+LiD & -- & $<$5 & \xmark & \cmark & \xmark \\
CrownViM \cite{crownvim} & Polygon & RGB & -- & 10 & \omark & \omark & \omark \\
BARE$^{\ddagger}$ \cite{bare} & Polygon & RGB & -- & 10 & \omark & \omark & \omark \\
\midrule
\multicolumn{8}{l}{\textit{Detection only -- no masks produced}} \\
DeepForest \cite{weinstein2020} & -- & RGB & -- & 10 & \xmark & \cmark & \xmark \\
SelvaBox \cite{selvabox} & -- & RGB & -- & 1--3 & \cmark & \cmark & \xmark \\
\midrule
\multicolumn{8}{l}{\textit{Foundation-model prompting}} \\
Box-prompt SAM \cite{teng2025} & Box & RGB+DSM & SAM & $<$5 & \xmark & \cmark & \xmark \\
LiDAR-prompt SAM \cite{sam2lidar} & Box & RGB+LiD & SAM\,2 & 10 & \xmark & \cmark & \xmark \\
SelvaMask \cite{selvamask} & Polygon & RGB & VFM & 1--4 & \cmark & \cmark & \xmark \\
\midrule
\multicolumn{8}{l}{\textit{Reduced supervision}} \\
ALS pseudo-labels \cite{alspseudo} & None$^{\dagger}$ & MS+ALS & SAM\,2 & -- & \xmark & \cmark & \xmark \\
FG-TreeSeg \cite{fgtreeseg} & Semantic & RGB & Cellpose & 10 & \xmark & \cmark & \xmark \\
\midrule
\rowcolor{black!8}
\textbf{LACE (ours)} & \textbf{Box} & \textbf{RGB} & \textbf{--} & \textbf{10}
 & \omark & \cmark & \cmark \\
\bottomrule
\end{tabular}

\vspace{3pt}
\begin{minipage}{0.95\textwidth}
\footnotesize
Tr, tropical; Te, temperate; Sv, savanna/dryland. $\bullet$ evaluated and
reported separately; $\circ$ present in an aggregate benchmark, not reported
as a stratum; $\times$ not evaluated. \emph{2nd net}: additional pretrained
network required at inference beyond the shared backbone. LiD, LiDAR;
MS, multispectral. $^{\dagger}$No manual instance annotation, but requires
co-registered ALS. $^{\ddagger}$Semantic segmentation only.
\end{minipage}
\end{table}
\subsubsection{Box to Mask}

- Box2Mask (TPAMI 2024, 42.4 mask AP COCO with Swin-L, "on par with fully mask-supervised")
- BoxSnake
- BoxInst
- DiscoBox.

\section{Materials and Methods}

% Materials and Methods should be described with sufficient details to allow others to replicate and build on published results. Please note that publication of your manuscript implicates that you must make all materials, data, computer code, and protocols associated with the publication available to readers. Please disclose at the submission stage any restrictions on the availability of materials or information. New methods and protocols should be described in detail while well-established methods can be briefly described and appropriately cited.

% In this section, where applicable, authors are required to disclose details of how gen-erative artificial intelligence (GenAI) has been used in this paper (e.g., to generate text, data, or graphics, or to assist in study design, data collection, analysis, or interpretation). The use of GenAI for superficial text editing (e.g., grammar, spelling, punctuation, and formatting) does not need to be declared.

\subsection{Datasets}

\paragraph{Restor OAM-TCD}

A significant recent contribution from the Restor Foundation in Zurich, Switzerland, has been the OAM-TCD dataset which uses OpenAerialMap (OAM, cite) RGB imagery and provides 5072 2048x2048px images at 0.1m/px resolution with associated human-labeled instance masks for over 280k individual and 56k groups of trees (groups are hereafter referred to as "canopy"). In our work, the OAM-TCD dataset has been an invaluable benchmark to assess the comparative performance of our LACE model in an apples-to-apples comparison against other SOTA models, due to the inclusion of not only bounding boxes but also precise polygon masks, as well as its broad range of ecological biomes to capture the diversity and morphology of trees in different terrestrial biomes including both urban and natural environments (Veitch-Michaelis, 2024). This multi-region, multi-biome distribution of annotated aerial RGB data allows for a much more diverse training set which aids cross-biome transfer, and also contains savannah \& arid dryland landscapes which are scarce in existing annotated datasets. 

Of the 5072 images, 439 are reserved as a holdout benchmark test set, allowing for a global multi-biome comparison between models. The 439 image test set contains 25,705 ITC annotations and 366 images within the set contain $\geq$1 crown i.e. 73 of the images have 0 ITCs. The same test set also contains 6,782 canopy annotations, which are groups of trees in a wide range of real-world sizes. 364 of 439 tiles contain $\geq$1 canopy region. Annotators were told to use the canopy class when it was not possible to reliably delineate individual trees, including when it was "not obvious whether a tree was an individual or multiple". 342 images contain at least one of each class and 51 tiles are fully empty.

% Research manuscripts reporting large datasets that are deposited in a publicly available database should specify where the data have been deposited and provide the relevant accession numbers. If the accession numbers have not yet been obtained at the time of submission, please state that they will be provided during review. They must be provided prior to publication.


\paragraph{NEON}TreeEvaluation

The U.S. National Science Foundation's National Ecological Observatory Network (NEON) publicly available aerial imagery dataset is a cornerstone of the RS research, comprised of a 22 separate locations across the continental United States and including a variety of different biomes. The NEONTreeEvaluation dataset (Weinstein et al., 2020) was pretrained on all 22 sites, using an unsupervised LiDAR based algorithm (Silva et al., 2016) to generate millions of moderate quality annotations for model pretraining, with a further 10,000 manual annotations of RGB imagery from six sites. The quality and scale of the NEON dataset as well as its fundamental role in the development of the DeepForest model was a key reason to build our proprietary datasets and CanopyAI to natively support the NEON resolution of 0.1m/px.

Data preprocessing: We pinned NeonTreeEvaluation @1.8.0 and built the evaluation GT by intersecting annotation XMLs with RGB tiles. This resulted in 194 scored tiles / 6,634 boxes / 22 sites. We downloaded the hand-annotated RGB training tiles and cropped them into 1,734 × 400px patches (22,831 boxes) restricted to genuinely-annotated regions only.  Two tiles were dropped due to low annotation coverage (SJER 41\% / TOOL 0\% coverage).  and capping empty negatives, with a leak-free tile-level train/val split.

Both are RGB-only (we did not use the dataset's LiDAR/CHM/HSI data at any time), at native 0.1 m/px, and the 400px patch geometry matches the 400px evaluation tiles so training and test see the same scale. Features are then extracted as padded-to-512 -> 32×32 DINOv3 grids (layers=21–24, parity-gated)%fix this last line

\subsection{LACE}
\label{sec:LACE}

The method converts a frozen self-supervised encoder into instance masks for individual tree crowns using bounding-box supervision only. It has three stages: a dense feature grid obtained by interlacing four offset passes of a frozen DINOv3-web ViT (Section~\ref{sec:features}); an anchor-free detector trained on bounding boxes (Section~\ref{sec:detector}); and a training-free box-to-mask module that converts each predicted box into a mask by foreground/background discrimination in a whitened feature space (Section~\ref{sec:masker}). No crown polygon is read at any stage of encoding, detection or mask fitting. The only mask labels used anywhere in the pipeline are the validation images on which the three inference scalars are selected (Section~\ref{sec:inference}). %wondering if perhaps this is distracting - alpha/kappa can maybe just be fixed one off..?

\subsubsection{Dense features by phase interlacing}
\label{sec:features}

A ViT with patch size $p$ emits one embedding per $p \times p$ patch, giving a feature grid $p$ times coarser than the image. For crowns whose diameter is a small multiple of $p$ this is the binding resolution constraint. Bilinear upsampling of that grid increases the sampling rate but introduces no new measurements.

Rather than upsample, we instead evaluate the frozen encoder four times per tile, shifting the input by $(\delta_y, \delta_x) \in \{0, p/2\}^2$, and interlace the four grids, which we found produced higher performance (Section~\ref{sec:ablation}). With $p = 16$ this yields a real $8$~px lattice. Writing $F^{(\delta_y,\delta_x)}$ for the grid from one offset, the assembled Feature tensor $\widetilde F$ is
\begin{equation}
\widetilde F\!\left[:,\; \tfrac{\delta_y}{8} \!:\! 2 \!:\! ,\; \tfrac{\delta_x}{8} \!:\! 2 \!:\! \right] \;=\; F^{(\delta_y,\delta_x)} ,
\label{eq:interlace}
\end{equation}

%fix this equation ^
so that cell $(Y, X)$ of $\widetilde F$ carries the embedding of the patch centred at pixel $(8X + 8,\, 8Y + 8)$. Each cell is therefore a genuine encoder output at its own location with a virtual 4~px border included, rather than an interpolation of its neighbours, and the union of the four phases' patch centres is exactly the $8$~px lattice.

Encoding uses DINOv3 ViT-L/16 with weights frozen throughout. We explored Domain-Adaptive Fine Tuning but found no material improvement upon the frozen weights (~\ref{sec:ablation}). We take final layer ( layer~24) activations ($1024$ dimensions), giving $\widetilde F \in \mathbb{R}^{1024 \times 256 \times 256}$ for a $2048$~px tile. Registration is verified independently: a synthetic blob is localised to within one cell across $200$ positions, and interlaced cells agree with bilinearly interpolated ones only at $\cos \approx 0.955$, confirming that the four passes contribute distinct evidence. %check this cos 0.955 figure source

\subsubsection{Detection}
\label{sec:detector}

Detection uses an anchor-free centre-based head operating directly on $\widetilde F$ at stride $8$. A $1\times1$ stem projects $1024$ channels to width $256$, followed by three $3\times3$ convolutional blocks with group normalisation and ReLU, and two $1\times1$ output heads: a single-channel centre heatmap and a four-channel regression map carrying a local sub-cell offset and log box size. This is based on the CenterNet design (cite CenterNet and CornerNet Law \& Deng). The heatmap bias is initialised to $-2.19$ (\%check this \%) so the initial centre probability is $\approx 0.1$, and the size biases are initialised to a $20$~px crown at stride $8$.

Training minimises a focal loss on the centre heatmap plus smooth-$L_1$ losses on offset and size at annotated centres, the size term weighted by $0.1$. Because the interlaced grid is already at the target stride, the head contains no internal upsampling; the half-cell difference between the assembled cell centre $8X+8$ and the target cell centre $8X+4$ is a uniform translation absorbed by the offset head, whose bias is initialised accordingly. Centre targets are encoded as CenterNet Gaussians whose radius follows the standard IoU criterion at an overlap of $0.7$, clamped never to fall below one feature cell so that the smallest crowns retain a non-degenerate target. Boxes are decoded by local-maximum selection on the heatmap, retaining candidates scoring above $0.05$ up to a maximum of $600$ per tile. Average precision is computed over the full ranked list; where precision, recall and F1 are reported at a single operating point, that point is a detector score of $0.4$.

\subsubsection{Training-free box-to-mask conversion}
\label{sec:masker}

Given a box, the task remains to decide which of its cells belong to the crown. We treat this as foreground/background discrimination in a whitened feature space, with the model fitted once on training tiles using boxes only and then held fixed.

The module has four key components: a whitened feature space; a frozen background mixture; a set of foreground prototypes; and a spatial prior over relative position within the box. The last two are learned by an alternating EM-based fit, while the first two stay fixed.

\paragraph{Feature space}
LACE's core idea is to leverage the rich DINOv3-web embeddings as sufficiently descriptive representations of each patch to build a high-performing tree-crown instance segmentation model. We begin with $a$, a given cell's raw unmodified embedding, and PCA-whiten it: $a$ is centred, projected onto the leading $d = 128$ principal directions of the training-cell distribution, and rescaled so that each retained component has unit variance. A final $L_2$ normalisation then places the result on the unit sphere, giving the unit vector $z$.
\begin{equation}
\tilde z \;=\; \frac{(a - \mu)\,W}{s}, \qquad z \;=\; \frac{\tilde z}{\lVert \tilde z \rVert},
\label{eq:whiten}
\end{equation}
Here $a \in \mathbb{R}^{1024}$ is one cell of $\widetilde F$; the cell Mean $\mu \in \mathbb{R}^{1024}$, the Whitening matrix $W \in \mathbb{R}^{1024 \times d}$ whose columns are the retained principal directions, and the Scale vector $s \in \mathbb{R}^{d}$ of the corresponding singular values are all estimated once on training cells and then frozen. This whitening/spherizing equalises the variance of the retained directions so that no direction dominates by scale alone. The final normalisation then places every cell on the unit sphere, making all subsequent comparisons simple cosine similarities and allowing both foreground and background to be modelled by von Mises--Fisher components (cite) with a shared concentration $\kappa$.

\paragraph{Background mixture}
Cells are partitioned between categories: in-box cells, clear background cells, and a ring of cells immediately outside each box. We exclude cells that fall within GT boxes so as not to accidentally define valid tree cells as background. The ring cells share illumination, phenology and local context (e.g. ground shadow) with the crown and are therefore suitable hard negatives. A background mixture of $K_{bg}$ components $\{(w_j, G_j)\}_{j=1}^{K_{bg}}$ is fitted to the union of background and ring cells by spherical $k$-means and \emph{frozen}, giving the background log-likelihood $\ell_{bg}$ as follows:
\begin{equation}
\ell_{bg}(z) \;=\; \log \textstyle\sum_j w_j \exp\!\left(\kappa\, z^{\top} G_j\right)
\label{eq:bgll}
\end{equation}
Each background centroid $G_j$ is a unit vector on the same sphere as the cells, the Mixture weights $w_j$ sum to one, and $z^{\top} G_j$ is therefore just a cosine similarity. The log-sum-exp acts as a smooth maximum over the components, so $\ell_{bg}(z)$ is a single scalar per cell reading as how well that cell matches its best-matching background component -- neither a threshold nor an objective, but the evidence term that Equation~\eqref{eq:logratio} scores the foreground against.\footnote{Because $G_j$, $w_j$ and $\kappa$ are all fixed once the mixture is frozen, and $z$ depends only on the frozen whitening of Equation~\eqref{eq:whiten}, every input to Equation~\eqref{eq:bgll} is determined before fitting begins: $\ell_{bg}$ is therefore computed once per tile, cached, and reused unchanged by every iteration. Note also that the standard von Mises--Fisher normalising constant $C_d(\kappa)$ is omitted here and from the foreground term of Equation~\eqref{eq:logratio} alike; sharing a single Concentration parameter $\kappa$ between the two mixtures makes it identical on both sides, so it cancels in the difference and never has to be evaluated in $d = 128$ dimensions. Equations~\eqref{eq:bgll}, \eqref{eq:logratio} and~\eqref{eq:resp} are written in the conventional mixture form $\sum_j w_j \exp(\cdot)$ for readability; all three are evaluated in log-sum-exp form for numerical stability.}

\paragraph{Foreground mixture}
Foreground and background take the same von Mises--Fisher form with a shared Concentration $\kappa$, and differ in one respect only: the background weights $w_j$ are global and frozen, whereas the foreground weights are supplied per cell by the Spatial prior and re-estimated at every M-step.
Foreground is represented by $K_{fg}$ unit Prototypes $\{C_k\}_{k=1}^{K_{fg}}$, the exact counterpart of the Background centroids $G_j$, initialised by spherical $k$-means over in-box cells restricted to those atypical of the background, i.e.\ cells which have a value of $\ell_{bg}$ below the box's median $\ell_{bg}$. This is merely a simple heuristic idea to kick-start the fitting - without this restriction, the initialisation is seeded from in-box background elements such as soil and roads, which are numerous and not the intended crown foreground.

\paragraph{Spatial prior}
We observe that a crown typically occupies a roughly central and roughly circular region of a reasonably tight GT box. We encode this as a spatial prior $\rho$ over the crown's \emph{relative} position within the box, to push the model towards producing reasonable tree crown-like prediction masks. In particular, we aim to increase LACE's robustness to dense-canopy and lush undergrowth situations, where the box-to-mask module can unintentionally produce masks that are not tree crown-like. The box is divided into an $N \times N$ grid of bins in normalised coordinates $(u, v) \in [0,1]^2$. Because those coordinates are a \emph{fraction} of the box, $\rho$ is by construction invariant to crown size, tile size and ground sample distance. Priors are maintained separately for three crown-size bands, whose edges are the terciles of the training box-size distribution, since at a fixed feature stride the number of cells spanning a crown depends on its size. Taken together, $\rho$ therefore carries one value for each (size band, prototype, $u$ bin, $v$ bin) combination: it is a joint prior over where in the box a cell sits and which prototype is expected there. We write
$\hat\rho_k = \rho_k / \sum_{k'} \rho_{k'}$
for the prior at a bin normalised across prototypes, so that $\hat\rho$ is a distribution over the $K_{fg}$ prototypes at that bin while the unnormalised total $\sum_k \rho_k$ remains the probability that the bin is crown at all. Equations~\eqref{eq:logratio} and~\eqref{eq:posterior} draw on these two quantities separately. $\rho$ is initialised uniform, so all prototypes begin with identical spatial profiles and differentiate only through Equation~\eqref{eq:priorupdate}.


Two constraints are applied to $\rho$ at every M-step. The first pins its mean per-bin foreground mass -- summed over prototypes and averaged over the $N^2$ bins -- to $\pi/4 \approx 0.785$, the area fraction of an ellipse inscribed in its bounding box and hence independent of box aspect ratio. The second caps any single bin at $0.95$, so that no relative position is ever treated as certainly foreground. Pinning the mass to a constant derived from geometry rather than fitted is what keeps the prior label-free, and it means the fit decides only \emph{where} the foreground mass sits, never \emph{how much} of it there is: without the constraint $\rho$ inflates towards unity, stops discriminating position, and the $\alpha$ term of Equation~\eqref{eq:posterior} degenerates into an additive constant.\footnote{The two constraints are not simultaneously satisfiable. On the deployed fit the rescaling is partly undone by the per-bin cap, leaving a realised mean per-bin mass of $0.690$ against the nominal $0.785$.}

\paragraph{Fitting}
Once we have these objects -- the whitening $(\mu, W, s)$, the frozen background mixture $\{(w_j, G_j)\}$, the foreground Prototypes $\{C_k\}$, and the Spatial prior $\rho$ -- we can run the EM fitting algorithm, which alternates a per-box E-step with a global M-step, set out in full in Algorithm~\ref{alg:fit}. Only $\{C_k\}$ and $\rho$ are learned here, the rest are estimated once and held fixed. 

\medskip\noindent\emph{E-step.} For the cells of one box, with $\rho_k$ the Spatial prior evaluated at each cell's bin for prototype $k$ and $\sum_k \rho_k$ the total foreground prior at that bin, the posterior follows in three steps: an appearance \emph{log-ratio}~\eqref{eq:logratio}, its within-box \emph{recentring}~\eqref{eq:recentre}, and the \emph{posterior} itself~\eqref{eq:posterior}.
\begin{align}
A(z) &\;=\; \log \textstyle\sum_k \hat\rho_k \exp\!\left(\kappa\, z^{\top} C_k\right) \;-\; \ell_{bg}(z),
\label{eq:logratio}\\[2pt]
\bar A(z) &\;=\; A(z) - \operatorname{mean}_{z' \in \mathrm{box}} A(z'),
\label{eq:recentre}\\[2pt]
p_{fg}(z) &\;=\; \sigma\!\left(\bar A(z) \;+\; \alpha \, \mathrm{logit}\!\left(\textstyle\sum_k \rho_k\right)\right).
\label{eq:posterior}
\end{align}
Equation~\eqref{eq:logratio} is a log-likelihood ratio. Its first term is the foreground log-likelihood, built from the Prototypes $C_k$ exactly as Equation~\eqref{eq:bgll} is built from the Background centroids $G_j$; subtracting $\ell_{bg}$ gives $A(z) = \log\!\big(p(z \mid \text{crown}) / p(z \mid \text{background})\big)$. A cell must therefore look crown-like \emph{and} unlike background to score highly. This is the first of two filters that in-box soil/ground has to pass, and it removes the soil/ground that resembles the background mixture outright.

Equation~\eqref{eq:recentre} is the second filter, and the step that makes the model work on frozen ViT features. Problematically, attention over the surrounding tile inflates the absolute foreground score of \emph{every} cell inside a box, so $A$ is confounded by context and its absolute level is not sufficiently informative: e.g. soil surrounded by canopy can still score positively, and therefore survive the first filter. However, recentering isolates $A$ \emph{within} a box, asking not whether a cell is crown-like but whether it is \emph{more} crown-like than the rest of its own box. The behaviour that follows is what we wanted to produce: a heterogeneous box containing crown and soil has a wide spread of $\bar A$ and is carved appropriately, whereas a box in uniformly closed canopy has a narrow spread and stays closer to filled, which is generally correct.

Equation~\eqref{eq:posterior}, the posterior, adds the only box-dependent term, and it is worth noting that the Spatial prior enters the two equations in different roles. In Equation~\eqref{eq:logratio} it appears only as the normalised weight $\hat\rho_k$, which sums to one over $k$ and therefore selects \emph{which} prototype applies at a given relative position while discarding \emph{how much} total foreground prior there is. That discarded total re-enters on its own in Equation~\eqref{eq:posterior}, where the Prior weight $\alpha$ scales it. The appearance term $\bar A$ is thus entirely box-independent, and Equation~\eqref{eq:posterior} contains the only point at which box geometry influences the posterior. Equation~\eqref{eq:posterior} gives the probability that a cell is crown, but the M-step also needs to know \emph{which} prototype accounts for it. The E-step therefore returns a second quantity, the Responsibilities $r_{nk}$\:
\begin{equation}
r_{nk} \;=\; p_{fg}(z_n)\;
\frac{\rho_k \exp\!\left(\kappa\, z_n^{\top} C_k\right)}
     {\sum_{k'} \rho_{k'} \exp\!\left(\kappa\, z_n^{\top} C_{k'}\right)},
\label{eq:resp}
\end{equation}
which distribute each cell's foreground probability over the $K_{fg}$ prototypes in proportion to the weighted components of Equation~\eqref{eq:logratio}, so that $\sum_k r_{nk} = p_{fg}(z_n)$. (Note that equation~\eqref{eq:resp} is written with $\rho_k$ rather than $\hat\rho_k$ because the normalisation appears in both numerator and denominator and therefore cancels.) Equations~\eqref{eq:logratio}--\eqref{eq:resp} are evaluated for all the cells of one box at a time and the results accumulated across all training boxes.

\medskip\noindent\emph{M-step.} The prototypes are updated in two steps. The first is the ordinary generative pull: prototype $k$ is drawn towards the responsibility-weighted mean of all cells over all boxes,
\begin{equation}
P_k \;=\; \frac{\sum_n r_{nk}\, z_n}{\lVert \sum_n r_{nk}\, z_n \rVert},
\label{eq:protomean}
\end{equation}
which expresses the idea ``recurs across boxes''. The second move repels the prototype from the outside-box signature it most resembles. The outside-box cells (clear background and ring) are softmax-assigned to the current prototypes at the same Concentration $\kappa$; the resulting per-prototype Negative centroid $\bar N_k$ is the normalised weighted mean of the cells assigned to prototype $k$, and the prototype becomes
\begin{equation}
C_k \;=\; \frac{P_k - \beta\,\bar N_k}{\lVert P_k - \beta\,\bar N_k \rVert}.
\label{eq:contrastive}
\end{equation}
A soil-like prototype's nearest negatives are themselves soil, so the subtraction tears it away from soil: commonality is defined as agreement-across-boxes \emph{minus} outside-box signature, not merely whatever fills the box. The Contrastive weight $\beta$ controls the strength of that repulsion; $\beta = 0$ recovers the purely generative update of Equation~\eqref{eq:protomean}, whose cost is reported in Section~\ref{sec:ablation}, and we use $\beta = 0.5$ throughout.

The Spatial prior is re-estimated from the same responsibilities. Writing $b(n)$ for the bin of cell $n$ and $N_b$ for the number of cells falling in bin $b$,
\begin{equation}
\tilde\rho_k(b) \;=\; \frac{1}{N_b} \sum_{n\,:\,b(n) = b} r_{nk},
\label{eq:priorupdate}
\end{equation}
accumulated separately per size band, after which $\tilde\rho$ is rescaled to the $\pi/4$ mean per-bin mass and capped at $0.95$ per bin as described above to give $\rho$. Finally, after a burn-in of five iterations, any prototype whose share $\sum_{n,b} r_{nk}$ of the total foreground responsibility falls below $0.25/K_{fg}$ of the uniform share is deleted, and $K_{fg}$ decreases accordingly.

Fitting is deterministic given the seed. The procedure is not gradient-based, reads no crown polygons, and produces a compact set of parameters $\{\mu, W, s, C, G, w, \rho\}$.

\begin{algorithm}[H]
\caption{Box-only masker fit. Inputs: training tiles, their boxes, and the frozen encoder. No crown polygon is read.\label{alg:fit}}
\begin{algorithmic}[1]
\State \textbf{Whitening.} Estimate $\mu$, $W$, $s$ on all training cells; map every cell by Equation~\eqref{eq:whiten}.
\State \textbf{Background.} Partition cells into in-box, ring and clear background by box geometry; fit $\{(w_j, G_j)\}$ on ring $\cup$ clear by spherical $k$-means; \textbf{freeze}. Precompute $\ell_{bg}$ once (Equation~\eqref{eq:bgll}) -- it never changes.
\State \textbf{Initialise $C$.} Spherical $k$-means over in-box cells with $\ell_{bg}$ below its median, giving $K_{fg}$ prototypes.
\State \textbf{Initialise $\rho$.} Uniform, at $\pi/4 \,/\, K_{fg}$ in every bin of every size band.
\For{$t = 1$ \textbf{to} $T = 30$}
  \For{each training box}                                \Comment{E-step, per box}
    \State Evaluate $A$, $\bar A$, $p_{fg}$, $r_{nk}$ by Equations~\eqref{eq:logratio}--\eqref{eq:resp}.
    \State Accumulate $\sum_n r_{nk} z_n$ and the per-bin sums of $r_{nk}$.
  \EndFor
  \State $P_k \gets$ Equation~\eqref{eq:protomean}; \; $C_k \gets$ Equation~\eqref{eq:contrastive}. \Comment{M-step, global}
  \State $\rho \gets$ Equation~\eqref{eq:priorupdate}, then rescale to $\pi/4$ and cap at $0.95$.
  \If{$t > 5$} delete prototypes below $0.25/K_{fg}$ of the uniform responsibility share. \EndIf
\EndFor
\State \Return $\{\mu, W, s, C, G, w, \rho\}$.
\end{algorithmic}
\end{algorithm}

Two important properties of this loop: the background mixture is frozen after step~2, so $\ell_{bg}$ is a constant of the fit rather than a quantity being jointly estimated. And the procedure is really EM-style rather than EM proper: the within-box recentring of Equation~\eqref{eq:recentre}, the contrastive term of Equation~\eqref{eq:contrastive} and the projection of $\rho$ onto the mass and cap constraints each break the correspondence with a joint likelihood, so no objective is monotonically ascended and there is no convergence test -- $T$ is fixed at $30$ iterations.

\subsubsection{Inference}
\label{sec:inference}

For each predicted box, cells within the box are projected by Equation~\eqref{eq:whiten}, the matching size band is selected, and posterior $p_{fg}$ is evaluated by Equations~\eqref{eq:logratio}--\eqref{eq:posterior}. Cells above a Mask threshold $\tau = 0.25$ form the predicted tree crown mask, which is then rasterised at the evaluation resolution. Only cells lying within the box are evaluated, and the rasterised mask is intersected with the box, so a mask is contained in its box by construction rather than by penalty.

Three scalars are exposed at inference, and they are the only quantities in the pipeline selected against mask labels. The Prior weight $\alpha$ scales the box-derived spatial term in Equation~\eqref{eq:posterior}: values $\leq 1$ relax reliance on a \emph{predicted} box, which is not expected to always be perfectly precise, so that the appearance term dominates. (Because the spatial prior is defined in normalised box coordinates, $\alpha$ is invariant to crown size and GSD -- what it does track is the precision of the boxes it is applied to). The concentration $\kappa$ is raised from the value of $10$ used during fitting to $16$ at inference, in Equations~\eqref{eq:bgll} and~\eqref{eq:logratio} alike; since $\bar A$ is zero-mean within each box, this scales the spread of the appearance term and thereby sharpens it. $10$ was fixed a priori, and $16$ is simply the best of $\{10, 13, 16, 20\}$ on validation, a range over which mask AP$_{50}$ varies by under $0.003$. The robust finding is directional -- any $\kappa > 10$ outperforms $\kappa = 10$ at every $\alpha$. 

The Mask threshold $\tau = 0.25$ is the posterior cut that defines the mask, lowered from $0.5$ so that small crowns under-covered by imprecise predicted boxes are recovered. All three are selected on held-out validation images (108/900) under a rule fixed in advance, and none requires refitting. We therefore describe the method as trained without mask supervision but broadly calibrated for tree crown mask generation on a small held-out set of mask labels, rather than as entirely label-free. We take these to be reasonable 'tree crown' hyperparameters and don't attempt to optimize them for use on other datasets such as NEONTreeEvaluation.


\subsubsection{Implementation}
\label{sec:implementation}

\begin{table}[H]
\caption{Configuration. Encoder and masker-fitting hyperparameters were fixed a priori;
the three inference scalars were selected on validation (Section~\ref{sec:inference});
detector optimisation follows standard practice for anchor-free centre-based
heads.\label{tab:config}}
\newcolumntype{K}{>{\hsize=0.42\hsize\raggedright\arraybackslash}X}
\newcolumntype{V}{>{\hsize=1.58\hsize\raggedright\arraybackslash}X}
\renewcommand{\arraystretch}{1.25}
\begin{tabularx}{\textwidth}{KV}
\toprule
\textbf{Component} & \textbf{Setting} \\
\midrule
Encoder & DINOv3 ViT-L/16, frozen, layer 24 only, $1024$ dimensions \\
\cmidrule{1-2}
Feature grid & 4 offsets $(\delta_y,\delta_x) \in \{0,8\}^2$, interlaced to $8$~px; $256\times256$ cells per $2048$~px tile \\
\cmidrule{1-2}
Detector & $1\times1$ stem to width $256$; $3\times$ $3\times3$ blocks, group norm, ReLU; heatmap and $4$-channel regression heads; stride $8$ \\
\cmidrule{1-2}
Target encoding & CenterNet Gaussians, grid $256$, stride $8$, radius at IoU overlap $0.7$, floor $\geq 1$ cell \\
\cmidrule{1-2}
Detector training & Adam, lr $10^{-3}$, weight decay $10^{-4}$, cosine schedule, batch $3$, max $40$ epochs, early stopping on validation box AP$_{50}$ (minimum $12$ epochs, patience $2$), best-on-validation checkpoint \\
\cmidrule{1-2}
Decoding & local maxima, score $> 0.05$, top $600$ per tile; operating point $0.4$ for single-threshold metrics \\
\midrule
Masker -- feature space & $120$ training tiles; PCA $d = 128$ \\
\cmidrule{1-2}
Masker -- mixtures & $K_{fg} = 16$ foreground prototypes; $K_{bg} = 12$ background components; concentration $\kappa = 10$; contrastive weight $\beta = 0.5$ \\
\cmidrule{1-2}
Masker -- spatial prior & $N = 8$ bins per axis; $3$ size bands at the training terciles; mean per-bin foreground mass $\pi/4$; per-bin cap $0.95$ \\
\cmidrule{1-2}
Masker -- EM schedule & $30$ iterations, no convergence test; prototype pruning after $5$ iterations below $0.25/K_{fg}$ of foreground responsibility \\
\midrule
Inference & prior weight $\alpha = 0.3$; concentration $\kappa = 16$ (fitted at $10$); mask threshold $\tau = 0.25$; masks rasterised at $512$~px per tile \\
\bottomrule
\end{tabularx}
\end{table}

The EM fit is CPU-only and completes in minutes. It is performed once, and the resulting masker is shared across all detector seeds, so the seed-to-seed variation reported in Section~\ref{sec:ablation} is detector variation. %check this hasn't been superceded, since we're looking at multiseed EM


\section{Results}


%%%%%%%%%%%%%%%%%%%%%%%%%%%%%%%%%%%%%%%%%%


\paragraph{Tree Crown Instance Segmentation}

For TCIS, only OAM-TCD is a suitable end-to-end evaluation dataset since it contains instance masks which are required for test-time evaluation. We restrict our evaluation to OAM-TCD's Tree class (cat=1) and report results on two approaches for completeness. On approach in the literature (cite SELVAMASK), which we support, is to ignore all canopy annotations (cat=2) altogether. Unsurprisingly, the model's predictions often fall inside areas marked as canopy, which contain many closely grouped and entirely valid tree crowns, but we do not attribute any positive or negative score to these predictions and they are excluded from the evaluation metrics entirely (SELVABOX delete all pixels from canopy areas, we use COCO's iscrowd=1 to simply ignore). The rationale here is that the annotating of canopy into groups necessarily leads to subjective groupings (for example should a clump of 10 tree crowns be best delineated by 1 canopy of 10 crowns, or 2 of 5? We claim both are equally valid) and the mAP50 metric becomes corrupted since IoU>0.5 is based on agreement of delineations. However, we also report figures without this consideration i.e. with canopy included and predictions inside canopy areas counting negatively towards score as False Positives since the OAM-TCD authors' Github code does not appear to make any exclusion in their reported result of 0.432 mAP50.

\begin{table}[H]
\begin{threeparttable}
\caption{Instance segmentation on the OAM-TCD 439-image holdout. \emph{Canopy neutral} means
predictions in unlabelled canopy are not penalised--by our canopy-ignore rule, or by SelvaMask's
equivalent deletion of the canopy class. \emph{Canopy = FP} counts them as false positives. Each
value sits in the column matching the protocol it was measured under. \emph{Imgs} is the number of
OAM-TCD training images seen.}
\label{tab:oamtcd439}
\newcolumntype{M}{>{\hsize=1.42\hsize\raggedright\arraybackslash}X}
\newcolumntype{S}{>{\hsize=0.58\hsize\raggedright\arraybackslash}X}
\setlength{\tabcolsep}{4pt}
\begin{tabularx}{\textwidth}{@{}M S r cc cc c@{}}
\toprule
 & & & \multicolumn{2}{c}{\textbf{Canopy Neutral}} & \multicolumn{2}{c}{\textbf{Canopy = FP}} & \\
\cmidrule(lr){4-5} \cmidrule(lr){6-7}
\textbf{Method} & \textbf{Superv.} & \textbf{Imgs} &
\textbf{AP\boldmath$_{50}$} & \textbf{AP\boldmath$_{50:95}$} &
\textbf{AP\boldmath$_{50}$} & \textbf{AP\boldmath$_{50:95}$} &
\shortstack{\textbf{Box}\\\textbf{AP\boldmath$_{50}$}} \\
\midrule
Restor Mask R-CNN\tnote{a} & masks & 4169 & & & 0.432 & & \\
DetecTree2 (fine-tuned) & masks & 900 & 0.535 & 0.224 & 0.421 & 0.180 & 0.529 \\
\addlinespace[2pt]
SelvaBox $\rightarrow$ SAM 3 (both FT)\tnote{b} & box + ext + SAM & $\sim$3335
  & & 0.185 & & & \\
SelvaBox $\rightarrow$ SAM 3 (frozen SAM)\tnote{b} & box + SAM & $\sim$3335
  & & 0.122 & & & \\
\textbf{LACE (ours)} (3-seed) & box & 900
  & \textbf{0.630} & \textbf{0.257} & \textbf{0.485} & \textbf{0.202} & \textbf{0.603} \\
 & & & {\footnotesize$\pm$0.005} & {\footnotesize$\pm$0.001}
     & {\footnotesize$\pm$0.002} & {\footnotesize$\pm$0.002} & \\
\bottomrule
\end{tabularx}
\begin{tablenotes}[flushleft]\footnotesize
\item[a] The OAM-TCD paper's own Mask R-CNN under its own protocol, without canopy-ignore; we have
not re-scored the released checkpoint, so only the canopy axis is known to align.
\item[b] Our protocol: 439 whole $2048^2$ tiles, pooled
COCO-101pt mask AP, masks at $512^2$, top-$k$ 600 detections at score $\geq$0.05. SelvaMask protocol: same 439 source images, different and easier test set: canopy deleted
and its pixels blacked out, cut to $1024^2$ at 0.5 overlap, tiles left empty or $>$80\% black dropped
(2527 of 3951 subtiles survive), AP averaged per tile rather than pooled (maxDets 400).
\end{tablenotes}
\end{threeparttable}
\end{table}


Tables~\ref{tab:oamtcd439} and~\ref{tab:neon-detection} summarize our model's performance across mAP (COCO mean Average Precision) for OAM-TCD and area/pixel-based semantic segmentation metrics for NEONTreeEvaluation, following the conventions established by the dataset authors' original publications. We report evaluations alongside baseline model results for comparison. For simplicity, consistency, and replicability, we use seed values starting at 0 and incremented by 1 (i.e. 0,1,2 for 3 seed).Unless noted otherwise, results are all generated under identical scoring regimes - in some cases we use reported figures and specify source publications. In these cases, we comment below on any deviations in testing regime based on public Github repo sources. 

\begin{table}[t]
\centering
\caption{Individual tree-crown detection on the NeonTreeEvaluation benchmark
(RGB-only, IoU 0.4, 194 tiles; precision/recall macro-averaged over images). Metrics at
each model's best-$F_1$ operating point; \emph{max\,R} is the recall ceiling (score
threshold $\to 0$)}
\label{tab:neon-detection}
\begin{tabular}{lcccc}
\toprule
Method & P & R & $F_1$ & max\,R \\
\midrule
DeepForest, published (v1.8.0)\textsuperscript{a}  & 0.659 & 0.790 & 0.719 & -- \\
DeepForest, replicated (v2.1.0)\textsuperscript{b} & 0.745 & 0.709 & 0.726 & 0.765 \\
Ours (5-seed)& $0.727\pm0.009$ & $0.729\pm0.009$ & $\mathbf{0.728}\pm0.003$ & $0.849\pm0.003$ \\
\bottomrule
\end{tabular}

\vspace{2pt}
{\footnotesize\raggedright
\textsuperscript{a}Weinstein et al.\ (2021, \textit{PLoS Comput.\ Biol.}), Table~3
image-annotated RGB operating point; full PR-curve / recall ceiling not reported (``--'').\quad
\textsuperscript{b}\texttt{deepforest} v2.1.0 pretrained release model, evaluated by us on
the same 194 tiles (deterministic; no seed variance). All rows scored with the benchmark
authors' \texttt{deepforest.evaluate\_boxes}.\par}
\end{table}


\subsection{Ablation Studies}
\label{sec:ablation}

We ablate the two central contributions: the interlaced feature grid and the training-free box-to-mask module. Unless stated otherwise, ablations of the module hold the detector fixed, so every arm scores the identical set of predicted boxes and any difference is attributable to mask generation alone. Bands are the sample standard deviation ($\mathrm{ddof}=1$) over independent seeds.

\subsubsection{Feature Grid: Interlacing Versus Interpolation}

DINOv3 emits one embedding per $16\times16$ patch. Our encoder evaluates the frozen backbone at four spatial offsets and interlaces the results into a real $8$~px grid. The obvious alternative is to simply interpolate the native $16$~px grid to $8$~px. Table~\ref{tab:grid} isolates that choice on NeonTreeEvaluation: encoder, detector, targets, tiles and scorer are held identical, and the encode/decode path is byte-identical between arms (\texttt{TargetConfig(grid=64, stride=8)}), so the comparison is purely feature quality.

\begin{table}[H]
\caption{Interlaced versus interpolated $8$~px features on NeonTreeEvaluation (194 tiles, 5 seeds per arm, detection F1 at IoU $0.4$). Both arms share the same detector, targets and scorer. \label{tab:grid}}
\newcolumntype{C}{>{\centering\arraybackslash}X}
\begin{tabularx}{\textwidth}{lCCC}
\toprule
\textbf{Feature grid} & \textbf{F1@0.4} & \textbf{max Recall} & \textbf{NIWO site F1} \\
\midrule
Native $16$~px, interpolated to $8$~px & $0.689 \pm 0.013$ & $0.779 \pm 0.009$ & $0.427 \pm 0.010$ \\
Interlaced $4$-offset, real $8$~px     & $\mathbf{0.728 \pm 0.003}$ & $\mathbf{0.849 \pm 0.003}$ & $\mathbf{0.588 \pm 0.017}$ \\
\midrule
$\Delta$ & $+0.039$ & $+0.070$ & $+0.161$ \\
\bottomrule
\end{tabularx}
\noindent{\footnotesize{Interlacing is not equivalent to upsampling: real and interpolated cells agree only at $\cos \approx 0.955$. Registration was verified independently (blob test, $\le 1$ cell over 200 positions). The NIWO gain is roughly $10\times$ the seed standard deviation and holds precision.}}
\end{table}

The gain is largest at NIWO, the densest and smallest-crown site. We observe here and elsewhere that this is the context where the $16$~px grid most struggles, due to the 8px cell grid under-resolving small individual crowns. Interpolation adds resolution to the sampling lattice but no new evidence, whereas the four offsets add evidence.

\subsubsection{Box-to-Mask Module}

The module's hyperparameters were fixed a priori except for the contrastive repulsion strength $\beta$, which we ablate here, and the three inference-time scalars of Section~\ref{sec:inference}: the spatial-prior weight $\alpha$, the concentration $\kappa$ -- raised from its fit-time value of $10$ to $16$ -- and the mask threshold $\tau$. Table~\ref{tab:masker} ablates $\beta$, $\alpha$ and $\kappa$ on OAM-TCD; $\tau$ is held at $0.25$ throughout. The $\beta$ rows share detector seed 0 and therefore identical boxes (box AP50 $=0.605$ for both), making that row pair an exact paired comparison.

\begin{table}[H]
\caption{Box-to-mask module on the OAM-TCD 439-tile test split. The $\beta$ arms share identical predicted boxes; the lower block reports multi-seed bands.\label{tab:masker}}
\newcolumntype{C}{>{\centering\arraybackslash}X}
\begin{tabularx}{\textwidth}{lCC}
\toprule
\textbf{Configuration} & \textbf{Mask AP50} & \textbf{Mask AP50:95} \\
\midrule
\multicolumn{3}{l}{\emph{Objective variant (seed 0, identical boxes; see Table~\ref{tab:components})}} \\
$\beta=0$ (generative M-step, absolute E-step)   & $0.5790$ & $0.2003$ \\
$\beta = 0.5$ (contrastive M-step, recentered E-step)       & $\mathbf{0.6203}$ & $\mathbf{0.2436}$ \\
\midrule
\multicolumn{3}{l}{\emph{Inference scalars, selected on 108 held-out validation images}} \\
$\alpha = 1.0$, $\kappa = 10$ (3 seeds) & $0.6232 \pm 0.0026$ & $0.2404 \pm 0.0039$ \\
$\alpha = 0.3$, $\kappa = 16$ (3 seeds) & $\mathbf{0.6301 \pm 0.0054}$ & $\mathbf{0.2568 \pm 0.0005}$ \\
\bottomrule
\end{tabularx}
\noindent{\footnotesize{$\alpha$ and $\kappa$ were chosen by a rule pre-registered before any validation cell was read, and the test split was not used for their selection. Both rows use detector seeds $0$--$2$, so the comparison is paired and the gain from the knobs is $+0.0069$ mask AP50; an earlier five-seed vanilla band ($0.615 \pm 0.011$) overstated it at $+0.015$ by including two seeds absent from the knobbed arm. The validation surface is a broad plateau (validation AP50 $0.630$--$0.639$ across the $\alpha \times \kappa$ grid; these are validation figures on 108 images and are not comparable to the test bands above, which they resemble only by coincidence), so these values are not uniquely optimal; what is robust is directional -- $\alpha<1$ beats $\alpha=1$ at every $\kappa$, and $\kappa>10$ beats $\kappa=10$ at every $\alpha$. The values themselves are coarse: $\kappa=16$ is simply the best of $\{10,13,16,20\}$, a range over which AP50 moves by under $0.003$. The mask threshold $\tau=0.25$ is the third such scalar; the validation grid spanned $\tau \in \{0.25, 0.30\}$, and the earlier move from $0.5$ to $0.25$ predates the validation protocol. Because these scalars are fitted against validation mask labels, we describe them as selected on 108 held-out images rather than as label-free.}}
\end{table}

The two rows of that first block differ in two components: our implementation ties the E-step recentring flag to $\beta$, so the $\beta = 0$ arm also disables recentring. Table~\ref{tab:components} separates them on the validation split.

\begin{table}[H]
\caption{M-step repulsion against E-step recentring, $K_{fg}=16$, EM seed 0, detector seed 0, 108 validation tiles. Mean per-crown mask IoU over 7041 crowns for the oracle-box columns and 5428 box-matched true positives for the predicted-box columns, with mask AP50 on predicted boxes in parentheses.\label{tab:components}}
\newcolumntype{C}{>{\centering\arraybackslash}X}
\begin{tabularx}{\textwidth}{llCC}
\toprule
\textbf{M-step} & \textbf{E-step} & \textbf{GT Oracle boxes} & \textbf{Predicted boxes} \\
\midrule
generative ($\beta = 0$)     & absolute  & $0.7402$ ($0.968$) & $0.6717$ ($0.594$) \\
generative ($\beta = 0$)     & recentred & $0.7895$ ($0.969$) & $0.7188$ ($0.641$) \\
contrastive ($\beta = 0.5$)  & absolute  & $0.7758$ ($0.983$) & $0.7063$ ($0.622$) \\
contrastive ($\beta = 0.5$)  & recentred & $0.7885$ ($0.967$) & $0.7193$ ($0.639$) \\
\midrule
\multicolumn{2}{l}{\emph{repulsion alone}}   & $+0.0356$ & $+0.0346$ \\
\multicolumn{2}{l}{\emph{recentring alone}}  & $+0.0493$ & $+0.0471$ \\
\multicolumn{2}{l}{\emph{both}}              & $+0.0483$ & $+0.0476$ \\
\multicolumn{2}{l}{\emph{interaction (both $-$ sum of singles)}} & $\mathbf{-0.0366}$ & $\mathbf{-0.0341}$ \\
\bottomrule
\end{tabularx}
\noindent{\footnotesize{Effects are relative to the generative/absolute cell. The two mechanisms are substitutes rather than complements: either alone recovers roughly $+0.04$, and both together recover the same $+0.048$, giving a large negative interaction in both box regimes. The interaction is the departure from additivity: were the two independent, enabling both would give $+0.0849$ on oracle boxes, whereas it gives $+0.0483$. Recentring is the slightly larger term throughout, consistent with the ablation on our development set, where removing recentring drives the fraction of mask mass in the box corners from $0.37$ to $0.61$ against $0.37$ to $0.41$ for removing repulsion. Oracle-box AP50 is saturated near $0.97$ because recall is perfect by construction there, so mean crown IoU is the discriminating quantity.}}
\end{table}


\subsubsection{Foreground Mixture Contribution}

The module was designed as a mixture of $K_{fg}$ foreground prototypes, but with a contrastive term (i.e. the variant used throughout) the mixture collapses, albeit without significant cost to performance. The mechanism is worth setting out, since subtracting something from each prototype might be expected to separate them rather than merge them. Because the outside-box cells are softmax-assigned at a shared concentration, the per-prototype Negative centroids $\bar N_k$ of Equation~\eqref{eq:contrastive} are near-identical, so the update injects a common component $-\beta \bar N_k$ into every prototype. The differences $P_k - P_{k'}$ survive the subtraction unchanged; it is the renormalisation that converts them into progressively smaller angles, as each prototype comes to point substantially along $-\bar N_k$. The effect compounds across iterations -- as the prototypes align, the negatives are assigned ever more uniformly and $\bar N_k$ becomes more exactly common -- and it is opposed throughout by the generative pull of Equation~\eqref{eq:protomean}, which continues to draw each prototype back towards its own responsibility-weighted mean. Which of the two dominates is what $\beta$ controls, and the balance is graded rather than sharp: mean pairwise cosine rises smoothly with $\beta$ (Table~\ref{tab:collapse}), reaching $0.994$ at the deployed $\beta = 0.5$, where the surviving set spans an effective rank of $1.66$ out of $16$. Table~\ref{tab:collapse} shows the effect for increasing values of $\beta$. Component pruning tests responsibility \emph{share} and near-duplicate prototypes each retain a healthy and roughly equal share, so prototypes largely remain unpruned. Pruning does fire at intermediate $\beta$, where some prototypes starve, but not at $\beta = 0.5$, where all sixteen are more or less equally engaged.

\begin{table}[H]
\caption{Foreground prototype geometry and accuracy, E-step recentring held on throughout so that $\beta$ is the only variable. Accuracy is mean per-crown mask IoU on the 108 validation tiles at EM seed 0 and detector seed 0, matching the geometry columns. Effective rank is $\exp$ of the entropy of the normalised singular-value spectrum of the prototype matrix; it counts directions rather than rows. \label{tab:collapse}}
\newcolumntype{C}{>{\centering\arraybackslash}X}
\begin{tabularx}{\textwidth}{lCCCCCCC}
\toprule
 & & & & & \multicolumn{2}{c}{\textbf{Mean crown IoU}} & \\
\cmidrule(lr){6-7}
\textbf{Config.} & \textbf{$K_{fg}$ init} & \textbf{$K_{fg}$ final} & \textbf{Eff. rank} & \textbf{Pairwise $\cos$} & \textbf{Oracle} & \textbf{Pred.} & \textbf{Pred. AP50} \\
\midrule
$\beta = 0$    & 16 & 16 & $11.43$ & $0.234$ & $0.7895$ & $0.7188$ & $0.6406$ \\
$\beta = 0.1$  & 16 & 15 & $9.41$  & $0.334$ & $0.7900$ & $0.7194$ & $0.6404$ \\
$\beta = 0.25$ & 16 & 13 & $6.37$  & $0.525$ & $0.7901$ & $0.7199$ & $0.6376$ \\
$\beta = 0.5$  & 16 & 16 & $\mathbf{1.66}$ & $\mathbf{0.994}$ & $0.7885$ & $0.7193$ & $0.6392$ \\
\midrule
$\beta = 0$    & 2  & 2  & $2.00$  & $0.053$ & $0.7888$ & $0.7181$ & $0.6369$ \\
$\beta = 0.5$  & 2  & 2  & $1.30$  & $0.988$ & $0.7891$ & $0.7196$ & $0.6384$ \\
\bottomrule
\end{tabularx}
\noindent{\footnotesize{Crown IoU is flat throughout: effective rank falls sevenfold across the $K_{fg}=16$ block while mean crown IoU spans $0.0011$ on predicted boxes and $0.0016$ on oracle boxes, both within the masker-seed band. Predicted AP50 spans $0.0030$ over that block, above its own seed band of $\pm 0.0004$, but it is not monotone in $\beta$ and is lowest at $\beta = 0.25$, where crown IoU is highest; it therefore does not order the configurations by degree of collapse. At $\beta = 0.5$ both $K_{fg}$ settings collapse, and the two reach the same direction ($\cos(\bar{C}) = 0.9942$, exceeding the between-seed agreement within $K_{fg}=16$ itself, $0.9915$--$0.9922$ over three EM seeds).}}
\end{table}

Surprisingly, refitting at three EM seeds with the detector fixed, the IoU of collapsed and non-collapsed configurations are essentially indistinguishable on the $108$ validation tiles: mean crown (3-seed) IoU $0.7204 \pm 0.0011$ at $\beta = 0.5$ against $0.7190 \pm 0.0004$ at $\beta = 0$, while mask AP50 runs the other way, $0.6390 \pm 0.0002$ against $0.6405 \pm 0.0004$. The two metrics disagree on which configuration leads, suggesting no real effect either way. Across the full single-seed sweep of Table~\ref{tab:collapse}, effective rank falls by a factor of seven for a spread of $0.0011$ in crown IoU. Masker refitting is also far less variable than detector training: across three EM seeds with the detector fixed, mask AP50 varies by $\pm 0.0004$, against the $\pm 0.0054$ across detector seeds in Table~\ref{tab:masker}. Fixing the masker at a single seed therefore does not meaningfully narrow the reported band. 

One reading of the collapse result is that a single ``not-background'' direction tolerates imprecise boxes where a diverse mixture would admit background. Table~\ref{tab:collapse} tests this directly, since recentring is held on there and $\beta$ is the only variable: collapse is worth $+0.0005$ mean crown IoU on predicted boxes against $-0.0010$ on oracle boxes. Were it protective, the first would exceed the second; instead the two differ by $0.0015$, disagree in sign, and both sit inside the masker-seed band. Whatever makes $\beta = 0$ and $\beta = 0.5$ score alike, it is not that one of them copes better with box error. The same independence holds for the larger gain over the generative/absolute configuration, $+0.0483$ on oracle boxes against $+0.0476$ on predicted (Table~\ref{tab:components}); that gain is the joint two-component effect rather than a property of $\beta$.


The foreground mixture is therefore not load-bearing. Neither the number of components nor the repulsion strength that determines whether they remain distinct changes the result. This is not because crown appearance is homogeneous: crowns vary in species, architecture, illuminated aspect and size, and the prototypes are global, shared across every box, so that variety is precisely what a foreground mixture might be expected to carry. It is that the diverse mixture was fitted and bought nothing. At $\beta = 0$ the prototypes remain near-orthogonal across eleven effective directions (mean pairwise $\cos 0.234$), yet score within $0.0005$ on predicted boxes and $0.0010$ on oracle boxes of the single direction that survives at $\beta = 0.5$; a genuine two-component fit ($K_{fg} = 2$, effective rank $2.00$ of $2.00$) lands in the same place. What makes the additional directions redundant is the E-step: recentring scores each crown against its own box statistics rather than matching it against a global appearance model, so variation between crowns is absorbed locally and need not be stored in the prototypes. The capacity asymmetry the data supports is thus the opposite of the one the design assumed --- a rich background model and a minimal foreground one --- with the discrimination carried by within-box recentring scored against the frozen background mixture, whose twelve components remain in genuine use (weight entropy $2.465$ of a possible $2.485$~nats).

We report this as a property of the objective rather than a defect of the implementation. Prototype collapse is well documented in self-supervised learning with this family of encoders, where it is uniformly treated as a pathology to be prevented. Here it arises from an explicit shared-negative subtraction rather than an optimisation shortcut, and it is inconsequential: the collapsed and non-collapsed configurations are equally accurate. We are not aware of a prior report of collapse without cost.

\subsubsection{Boundary Quality and Behaviour at Strict IoU}

Table~\ref{tab:boundary} reports boundary F1 on matched true positives, which conditions on successful detection and so isolates mask quality from recall.

\begin{table}[H]
\caption{Boundary F1 on matched true positives (IoU $\ge 0.5$), OAM-TCD 439 tiles, scored on the $512$~px raster. A tolerance of $1$~px here corresponds to $4$~px in tile coordinates. \label{tab:boundary}}
\newcolumntype{C}{>{\centering\arraybackslash}X}
\begin{tabularx}{\textwidth}{lCCC}
\toprule
\textbf{Method} & \textbf{Supervision} & \textbf{BF@1px} & \textbf{BF@2px} \\
\midrule
Filled predicted box              & --         & $0.6576 \pm 0.0220$ & $0.8306$ \\
Ours (3 seeds)                    & boxes       & $0.7611 \pm 0.0020$ & $0.9019$ \\
DetecTree2, fine-tuned (1 seed)   & full masks  & $0.7666$            & $0.9082$ \\
\bottomrule
\end{tabularx}
\noindent{\footnotesize{Ours and the filled-box baseline are evaluated on identical matched pairs and are exactly paired; DetecTree2 matches its own detections (19{,}217 true positives against approximately 20{,}750) and that comparison is therefore approximate. Tolerances above $1$~px saturate ($\ge 0.90$ at $2$~px, $\ge 0.99$ at $5$~px) and are reported for completeness only.}}
\end{table}

Box supervision costs essentially nothing in boundary placement: we exceed a filled box by $0.104$ at the tightest measurable tolerance and sit within $0.006$ of a mask-supervised model. The cost appears instead at strict IoU. Our lead over DetecTree2 is $+0.096$ AP at IoU $0.50$ and decays monotonically, with the two methods crossing at IoU $0.75$. This is a resolution limit rather than a masker deficiency: with oracle boxes our mean per-crown mask IoU is $0.789$ (Table~\ref{tab:components}), and across matched true positives $39.2\%$ exceed IoU $0.75$ while only $20.6\%$ exceed $0.80$. The $8$~px feature grid bounds attainable IoU, so the method competes on foreground discrimination and spatial priors rather than on pixel-exact delineation.

\subsubsection{Dense canopy}

Foreground/background commonality within a box rests to some extent on the assumption that the box contains mostly one crown, meaning crowded trees with overlapping crowns could be a challenge.


The assumption is a claim about individual crowns, so the direct test is at the instance level. We call a crown \emph{touching} when its ground-truth box overlaps that of another crown, which holds for $39.3\%$ of the 25{,}705 test crowns. Touching and isolated crowns score almost identically: mean mask IoU $0.7093 \pm 0.0015$ against $0.7148 \pm 0.0019$, and they fail at statistically indistinguishable rates ($0.0250 \pm 0.0031$ versus $0.0252 \pm 0.0008$, where failure is a box-matched true positive whose mask IoU falls below $0.5$). Touching crowns are also substantially larger in the mean ($70.2$ versus $44.1$~px), which we theorize could partly offset neighbour interference.

We can examine tile-level canopy closure to give a (coarser) corroboration of the same conclusion (Table~\ref{tab:strata}).

\begin{table}[H]
\caption{Performance by canopy closure, 439 OAM-TCD test tiles, 3 seeds. AP is pooled within each bin against that bin's own ground-truth count. The mask-minus-box gap is the module's contribution once detection is accounted for.\label{tab:strata}}
\newcolumntype{C}{>{\centering\arraybackslash}X}
\begin{tabularx}{\textwidth}{lCCCC}
\toprule
\textbf{Canopy closure} & \textbf{Tiles} & \textbf{Mask AP50} & \textbf{Box AP50} & \textbf{Mask $-$ Box} \\
\midrule
$[0,\,0.10)$    & 196 & $0.5519$ & $0.5271$ & $+0.0247$ \\
$[0.10,\,0.25)$ &  95 & $0.6601$ & $0.6359$ & $+0.0242$ \\
$[0.25,\,0.50)$ &  86 & $0.6824$ & $0.6562$ & $+0.0261$ \\
$[0.50,\,1.00]$ &  62 & $0.6362$ & $0.6036$ & $+0.0326$ \\
\bottomrule
\end{tabularx}
\noindent{\footnotesize{Bins are percentage of tile covered by tree crown; tile counts sum to the 439 test tiles. Box AP50 is masker-invariant and is reported so the module's contribution can be read independently of detection.}}
\end{table}

Table~\ref{tab:strata} shows higher canopy coverage does not cause and decline in AP50 -- perhaps even slightly the opposite. Mask AP50 is lowest on the most open tiles ($0.5519$ in the $[0,\,0.10)$ bin) and higher wherever canopy is more closed, peaking at $0.6824$ in the $[0.25,\,0.50)$ bin. That variation is a detection effect rather than a mask one, since box AP50 moves with it bin for bin, and the mask-minus-box gap---the masker's own contribution---stays between $+0.024$ and $+0.033$ across the four bins and in fact rises slightly as canopy closes.


\section{Discussion}

%Authors should discuss the results and how they can be interpreted from the perspective of previous studies and of the working hypotheses. The findings and their implications should be discussed in the broadest context possible. Future research directions may also be highlighted.

%%%%%%%%%%%%%%%%%%%%%%%%%%%%%%%%%%%%%%%%%%

\subsection{TCIS}

By taking the $8$~px grid as the unit of crown delineation, LACE concedes an inbuilt disadvantage relative to pixel-based maskers: in its current form it is fundamentally limited to a coarser resolution than the imagery supports. This would be expected to cost most at high IoU thresholds --- AP$_{75}$, AP$_{90}$ --- where pixel-level precision dominates the metric, and comparatively little at AP$_{50}$, where placement matters more than boundary exactness. Section~\ref{sec:ablation} reports where the predicted decay is and is not observed.

\subsubsection{Spatial Prior}

Although the Spatial prior is maintained separately for three crown-size bands, in practice we observe very little difference between the three fitted priors. The stratification is therefore cheap insurance rather than a working component of the method, and a single pooled prior would likely perform comparably.



\subsection{Box-to-mask}




\subsection{OAM-TCD Benchmark}



Comment on existing canopy annotations

Helpfully, NEONTreeEvaluation contains many GT annotations within dense canopy areas, annotated individually allowing for precise individual instance segmentation evaluation even in canopy areas.  Existing datasets that try to provide instance segmentation labels very often will tend to attempt to draw definitive 'ground truth' on dense canopy where there arguably is simply not enough information in the image to make that determination definitively (Yi et al. 2023), Ball et al., 2023). There are a variety of potential solutions to this challenge, and naturally, there is a need for pragmatism in annotating a close approximation of the actual ground truth in a scenario where it is impossible to be certain. Nevertheless, we observe that it is characteristic of datasets that annotate individual tree crowns in dense canopy scenarios, either with rectangular bounding boxes or closed polygons/masks, to leave a large number of gaps between units. As an approximation and indeed for pixel-wise classification loss, this can be sufficient, but introduces some important nuance when it comes to inference-time test set benchmarks. 

Firstly, where GT annotations leave non-trivial gaps between annotation boundaries, a model that produces more inclusive,  justified boundaries can be penalized. Secondly, COCOtools standard mean average precision (mAP50) is based on Intersection of Union (IoU) and, for each ground truth (can be bounding box or polygon), algorithmically looks for predictions that overlap with the ground truth, ordered by confidence score (descending) - if one is found that exceeds 0.5 i.e. the area of intersection between prediction and ground truth (GT) is 50\% or greater than the union of the two, that prediction is counted as a True Positive (TP). In annotated dense canopy, this entails that the resulting benchmark contains some subjectivity i.e.  many predictions can be counted as False Positive (FP) which are defensible and possibly even more accurate than the ground truth. This is particularly an issue with RGB imagery that includes trees of varying scales i.e. smaller and larger, as a mismatched scale between prediction and GT leads to much higher FP rates since mAP50 is highly sensitive to scale mismatch for purely geometric reasons - a small prediction's overlap with a large GT is much less likely to cross the 50\% IoU threshold and equally a large prediction's overlap with a small GT. If one of the two is less than 50\% the area of the other, it will be counted as a FP regardless of placement. In contexts such as dense canopy where, for example, a single large crown could be one tree or two trees, the metric is sensitive to the arbitrary choice of division. The same is true of Precision, Recall, and F1 metrics derived from instance predictions (not including area-/pixel-based Precision, Recall, or F1 metrics which are calculated by total pixels rather than  individual bounding boxes or polygons).


\section{Conclusion}

% check consistent useage of 'weakly supervised'
In this study we present LACE, a domain-agnostic model for instance segmentation in weakly supervised Tree Crown tasks, where annotations are scarce. Our box-to-mask inference method, by leveraging the statistical signature of the DINOv3 latent composite, provides a SOTA solution to true Tree Crown Instance Segmentation (TCIS) without the need for time-consuming and/or expensive polygon masks. With this contribution, we hope to support the Remote Sensing community in the previously underexplored area of TCD and TCIS in Savanna \& Dryland biomes. We also believe the designs discussed here would be highly applicable and transferrable to other Computer Vision domains such as medical imaging where availability of annotated data is also a known challenge, since we did not rely on domain-specific priors to any significant degree.

The efficacy of each component within LACE has been validated through comprehensive ablation studies. %(make sure this is correct).
%Simultaneously, experimental results on three public remote sensing datasets substantiate that our method surpasses other state-of-the-art instance segmentation methods, as well as several additional SAM-based methods.



%%%%%%%%%%%%%%%%%%%%%%%%%%%%%%%%%%%%%%%%%%
%% optional
%\supplementary{The following supporting information can be downloaded at:  \linksupplementary{s1}, Figure S1: title; Table S1: title; Video S1: title.}

% Only for journal Methods and Protocols:
% If you wish to submit a video article, please do so with any other supplementary material.
% \supplementary{The following supporting information can be downloaded at: \linksupplementary{s1}, Figure S1: title; Table S1: title; Video S1: title. A supporting video article is available at doi: link.}

% Only used for preprtints:
% \supplementary{The following supporting information can be downloaded at the website of this paper posted on \href{https://www.preprints.org/}{Preprints.org}.}

% Only for journal Hardware:
% If you wish to submit a video article, please do so with any other supplementary material.
% \supplementary{The following supporting information can be downloaded at: \linksupplementary{s1}, Figure S1: title; Table S1: title; Video S1: title.\vspace{6pt}\\
%\begin{tabularx}{\textwidth}{lll}
%\toprule
%\textbf{Name} & \textbf{Type} & \textbf{Description} \\
%\midrule
%S1 & Python script (.py) & Script of python source code used in XX \\
%S2 & Text (.txt) & Script of modelling code used to make Figure X \\
%S3 & Text (.txt) & Raw data from experiment X \\
%S4 & Video (.mp4) & Video demonstrating the hardware in use \\
%... & ... & ... \\
%\bottomrule
%\end{tabularx}
%}

%%%%%%%%%%%%%%%%%%%%%%%%%%%%%%%%%%%%%%%%%%
\authorcontributions{For research articles with several authors, a short paragraph specifying their individual contributions must be provided. The following statements should be used ``Conceptualization, X.X. and Y.Y.; methodology, X.X.; software, X.X.; validation, X.X., Y.Y. and Z.Z.; formal analysis, X.X.; investigation, X.X.; resources, X.X.; data curation, X.X.; writing--original draft preparation, X.X.; writing--review and editing, X.X.; visualization, X.X.; supervision, X.X.; project administration, X.X.; funding acquisition, Y.Y. All authors have read and agreed to the published version of the manuscript.'', please turn to the  \href{http://img.mdpi.org/data/contributor-role-instruction.pdf}{CRediT taxonomy} for the term explanation. Authorship must be limited to those who have contributed substantially to the work~reported.}

\funding{Please add: ``This research received no external funding'' or ``This research was funded by NAME OF FUNDER grant number XXX.'' and  and ``The APC was funded by XXX''. Check carefully that the details given are accurate and use the standard spelling of funding agency names at \url{https://search.crossref.org/funding}, any errors may affect your future funding.}

\institutionalreview{In this section, you should add the Institutional Review Board Statement and approval number, if relevant to your study. You might choose to exclude this statement if the study did not require ethical approval. Please note that the Editorial Office might ask you for further information. Please add “The study was conducted in accordance with the Declaration of Helsinki, and approved by the Institutional Review Board (or Ethics Committee) of NAME OF INSTITUTE (protocol code XXX and date of approval).” for studies involving humans. OR “The animal study protocol was approved by the Institutional Review Board (or Ethics Committee) of NAME OF INSTITUTE (protocol code XXX and date of approval).” for studies involving animals. OR “Ethical review and approval were waived for this study due to REASON (please provide a detailed justification).” OR “Not applicable” for studies not involving humans or animals.}

\informedconsent{Any research article describing a study involving humans should contain this statement. Please add ``Informed consent was obtained from all subjects involved in the study.'' OR ``Patient consent was waived due to REASON (please provide a detailed justification).'' OR ``Not applicable'' for studies not involving humans. You might also choose to exclude this statement if the study did not involve humans.

Written informed consent for publication must be obtained from participating patients who can be identified (including by the patients themselves). Please state ``Written informed consent has been obtained from the patient(s) to publish this paper'' if applicable.}

\dataavailability{We encourage all authors of articles published in MDPI journals to share their research data. In this section, please provide details regarding where data supporting reported results can be found, including links to publicly archived datasets analyzed or generated during the study. Where no new data were created, or where data is unavailable due to privacy or ethical restrictions, a statement is still required. Suggested Data Availability Statements are available in section ``MDPI Research Data Policies'' at \url{https://www.mdpi.com/ethics}.} 

%\durcstatement{Current research is limited to the [please insert a specific academic field, e.g., XXX], which is beneficial [share benefits and/or primary use] and does not pose a threat to public health or national security. Authors acknowledge the dual-use potential of the research involving xxx and confirm that all necessary precautions have been taken to prevent potential misuse. As an ethical responsibility, authors strictly adhere to relevant national and international laws about DURC. Authors advocate for responsible deployment, ethical considerations, regulatory compliance, and transparent reporting to mitigate misuse risks and foster beneficial outcomes.}

\acknowledgments{In this section you can acknowledge any support given which is not covered by the author contribution or funding sections. This may include administrative and technical support, or donations in kind (e.g., materials used for experiments). Where GenAI has been used for purposes such as generating text, data, or graphics, or for study design, data collection, analysis, or interpretation of data, please add “During the preparation of this manuscript/study, the author(s) used [tool name, version information] for the purposes of [description of use]. The authors have reviewed and edited the output and take full responsibility for the content of this publication.”}

\conflictsofinterest{Declare conflicts of interest or state ``The authors declare no conflicts of interest.'' Authors must identify and declare any personal circumstances or interest that may be perceived as inappropriately influencing the representation or interpretation of reported research results. Any role of the funders in the design of the study; in the collection, analyses or interpretation of data; in the writing of the manuscript; or in the decision to publish the results must be declared in this section. If there is no role, please state ``The funders had no role in the design of the study; in the collection, analyses, or interpretation of data; in the writing of the manuscript; or in the decision to publish the results''.} 

%%%%%%%%%%%%%%%%%%%%%%%%%%%%%%%%%%%%%%%%%%
%% Optional

%% Only for journal Encyclopedia
%\entrylink{The Link to this entry published on the encyclopedia platform.}

\abbreviations{Abbreviations}{
The following abbreviations are used in this manuscript:
\\

\noindent 
\begin{tabular}{@{}ll}
MDPI & Multidisciplinary Digital Publishing Institute\\
DOAJ & Directory of open access journals\\
TLA & Three letter acronym\\
LD & Linear dichroism
\end{tabular}
}

%%%%%%%%%%%%%%%%%%%%%%%%%%%%%%%%%%%%%%%%%%
%% Optional
\appendixtitles{no} % Leave argument "no" if all appendix headings stay EMPTY (then no dot is printed after "Appendix A"). If the appendix sections contain a heading then change the argument to "yes".
\appendixstart
\appendix
\section[\appendixname~\thesection]{}
\subsection[\appendixname~\thesubsection]{}
The appendix is an optional section that can contain details and data supplemental to the main text--for example, explanations of experimental details that would disrupt the flow of the main text but nonetheless remain crucial to understanding and reproducing the research shown; figures of replicates for experiments of which representative data are shown in the main text can be added here if brief, or as Supplementary Data. Mathematical proofs of results not central to the paper can be added as an appendix.

\begin{table}[H] 
\caption{This is a table caption.\label{tab5}}
%\newcolumntype{C}{>{\centering\arraybackslash}X}
\begin{tabularx}{\textwidth}{CCC}
\toprule
\textbf{Title 1}	& \textbf{Title 2}	& \textbf{Title 3}\\
\midrule
Entry 1		& Data			& Data\\
Entry 2		& Data			& Data\\
\bottomrule
\end{tabularx}
\end{table}

\section[\appendixname~\thesection]{}
All appendix sections must be cited in the main text. In the appendices, Figures, Tables, etc. should be labeled, starting with ``A''--e.g., Figure A1, Figure A2, etc.

%%%%%%%%%%%%%%%%%%%%%%%%%%%%%%%%%%%%%%%%%%
%\isPreprints{}{% This command is only used for ``preprints''.
\begin{adjustwidth}{-\extralength}{0cm}
%} % If the paper is ``preprints'', please uncomment this parenthesis.
%\printendnotes[custom] % Un-comment to print a list of endnotes

\reftitle{References}

% Please provide either the correct journal abbreviation (e.g. according to the “List of Title Word Abbreviations” http://www.issn.org/services/online-services/access-to-the-ltwa/) or the full name of the journal.
% Citations and References in Supplementary files are permitted provided that they also appear in the reference list here. 

%=====================================
% References, variant A: external bibliography
%=====================================
% \bibliography{your_external_BibTeX_file}

%=====================================
% References, variant B: internal bibliography
%=====================================

% ACS format
\isAPAandChicago{}{%
\begin{thebibliography}{999}
% Reference 1
\bibitem[Author1(year)]{ref-journal}
Author~1, T. The title of the cited article. {\em Journal Abbreviation} {\bf 2008}, {\em 10}, 142--149.
% Reference 2
\bibitem[Author2(year)]{ref-book1}
Author~2, L. The title of the cited contribution. In {\em The Book Title}; Editor 1, F., Editor 2, A., Eds.; Publishing House: City, Country, 2007; pp. 32--58.
% Reference 3
\bibitem[Author1 and Author2 (year)]{ref-book2}
Author 1, A.; Author 2, B. \textit{Book Title}, 3rd ed.; Publisher: Publisher Location, Country, 2008; pp. 154--196.
% Reference 4
\bibitem[Author4(year)]{ref-unpublish}
Author 1, A.B.; Author 2, C. Title of Unpublished Work. \textit{Abbreviated Journal Name} year, \textit{phrase indicating stage of publication (submitted; accepted; in press)}.
% Reference 5
\bibitem[Author8(year)]{ref-url}
Title of Site. Available online: URL (accessed on Day Month Year).
% Reference 6
\bibitem[Author6(year)]{ref-proceeding}
Author 1, A.B.; Author 2, C.D.; Author 3, E.F. Title of presentation. In Proceedings of the Name of the Conference, Location of Conference, Country, Date of Conference (Day Month Year); Abstract Number (optional), Pagination (optional).
% Reference 7
\bibitem[Author7(year)]{ref-thesis}
Author 1, A.B. Title of Thesis. Level of Thesis, Degree-Granting University, Location of University, Date of Completion.
\end{thebibliography}
}

% Chicago format (Used for journal: arts, genealogy, histories, humanities, jintelligence, laws, literature, religions, risks, socsci)
\isChicagoStyle{%
\begin{thebibliography}{999}
% Reference 1
\bibitem[Aranceta-Bartrina(1999a)]{ref-journal}
Aranceta-Bartrina, Javier. 1999a. Title of the cited article. \textit{Journal Title} 6: 100--10.
% Reference 2
\bibitem[Aranceta-Bartrina(1999b)]{ref-book1}
Aranceta-Bartrina, Javier. 1999b. Title of the chapter. In \textit{Book Title}, 2nd ed. Edited by Editor 1 and Editor 2. Publication place: Publisher, vol. 3, pp. 54–96.
% Reference 3
\bibitem[Baranwal and Munteanu {[1921]}(1955)]{ref-book2}
Baranwal, Ajay K., and Costea Munteanu. 1955. \textit{Book Title}. Publication place: Publisher, pp. 154--96. First published 1921 (op-tional).
% Reference 4
\bibitem[Berry and Smith(1999)]{ref-thesis}
Berry, Evan, and Amy M. Smith. 1999. Title of Thesis. Level of Thesis, Degree-Granting University, City, Country. Identifi-cation information (if available).
% Reference 5
\bibitem[Cojocaru et al.(1999)]{ref-unpublish}
Cojocaru, Ludmila, Dragos Constatin Sanda, and Eun Kyeong Yun. 1999. Title of Unpublished Work. \textit{Journal Title}, phrase indicating stage of publication.
% Reference 6
\bibitem[Driver et al.(2000)]{ref-proceeding}
Driver, John P., Steffen Rohrs, and Sean Meighoo. 2000. Title of Presentation. In \textit{Title of the Collected Work} (if available). Paper presented at Name of the Conference, Location of Conference, Date of Conference.
% Reference 7
\bibitem[Harwood(2008)]{ref-url}
Harwood, John. 2008. Title of the cited article. Available online: URL (accessed on Day Month Year).
\end{thebibliography}
}{}

% APA format (Used for journal: admsci, behavsci, businesses, econometrics, economies, education, ejihpe, games, humans, ijfs, journalmedia, jrfm, languages, psycholint, publications, tourismhosp, youth)
\isAPAStyle{%
\begin{thebibliography}{999}
% Reference 1
\bibitem[\protect\citeauthoryear{Azikiwe \BBA\ Bello}{{2020a}}]{ref-journal}
Azikiwe, H., \& Bello, A. (2020a). Title of the cited article. \textit{Journal Title}, \textit{Volume}(Issue), 
Firstpage--Lastpage/Article Number.
% Reference 2
\bibitem[\protect\citeauthoryear{Azikiwe \BBA\ Bello}{{2020b}}]{ref-book1}
Azikiwe, H., \& Bello, A. (2020b). \textit{Book title}. Publisher Name.
% Reference 3
\bibitem[Davison(1623/2019)]{ref-book2}
Davison, T. E. (2019). Title of the book chapter. In A. A. Editor (Ed.), \textit{Title of the book: Subtitle} 
(pp. Firstpage--Lastpage). Publisher Name. (Original work published 1623) (Optional).
% Reference 4
\bibitem[Fistek et al.(2017)]{ref-proceeding}
Fistek, A., Jester, E., \& Sonnenberg, K. (2017, Month Day). Title of contribution [Type of contribution]. Conference Name, Conference City, Conference Country.
% Reference 5
\bibitem[Hutcheson(2012)]{ref-thesis}
Hutcheson, V. H. (2012). \textit{Title of the thesis} [XX Thesis, Name of Institution Awarding the Degree].
% Reference 6
\bibitem[Lippincott \& Poindexter(2019)]{ref-unpublish}
Lippincott, T., \& Poindexter, E. K. (2019). \textit{Title of the unpublished manuscript} [Unpublished manuscript/Manuscript in prepara-tion/Manuscript submitted for publication]. Department Name, Institution Name.
% Reference 7
\bibitem[Harwood(2008)]{ref-url}
Harwood, J. (2008). \textit{Title of the cited article}. Available online: URL (accessed on Day Month Year).
\end{thebibliography}
}{}

% If authors have biography, please use the format below
%\section*{Short Biography of Authors}
%\bio
%{\raisebox{-0.35cm}{\includegraphics[width=3.5cm,height=5.3cm,clip,keepaspectratio]{Definitions/author1.pdf}}}
%{\textbf{Firstname Lastname} Biography of first author}
%
%\bio
%{\raisebox{-0.35cm}{\includegraphics[width=3.5cm,height=5.3cm,clip,keepaspectratio]{Definitions/author2.jpg}}}
%{\textbf{Firstname Lastname} Biography of second author}

% For the MDPI journals use author-date citation, please follow the formatting guidelines on http://www.mdpi.com/authors/references
% To cite two works by the same author: \citeauthor{ref-journal-1a} (\citeyear{ref-journal-1a}, \citeyear{ref-journal-1b}). This produces: Whittaker (1967, 1975)
% To cite two works by the same author with specific pages: \citeauthor{ref-journal-3a} (\citeyear{ref-journal-3a}, p. 328; \citeyear{ref-journal-3b}, p.475). This produces: Wong (1999, p. 328; 2000, p. 475)

%%%%%%%%%%%%%%%%%%%%%%%%%%%%%%%%%%%%%%%%%%
%% for journal Sci
%\reviewreports{\\
%Reviewer 1 comments and authors’ response\\
%Reviewer 2 comments and authors’ response\\
%Reviewer 3 comments and authors’ response
%}
%%%%%%%%%%%%%%%%%%%%%%%%%%%%%%%%%%%%%%%%%%
\PublishersNote{}
%\isPreprints{}{% This command is only used for ``preprints''.
\end{adjustwidth}
%} % If the paper is ``preprints'', please uncomment this parenthesis.
\end{document}

